\documentclass[output=paper,colorlinks,citecolor=brown,
% hidelinks,
% showindex
]{langscibook}
\ChapterDOI{10.5281/zenodo.20285184}

%This is where you put the authors and their affiliations
\author{Cassandra Caragine\affiliation{University of Maryland}}

%Insert your title here
\title{Vowel alternations in {D}agaare number marking}  
\abstract{Dagaare (Mabia) has two numbers, singular and plural, but nouns form the singular and plural in a variety of ways depending on their declension class. In this paper, I provide novel data from Northern Dagaare to investigate the vowel alternations exhibited by three of the classes: -r marked singulars, -bɛ marked plurals, and -r(ɪ) marked plurals. Following previous work on Central Dagaare \citep{anttilaProsodic2009, anttilaMetrically2019}, I conclude that the vowel alternations in the -r singulars are in fact true suffixes. However, \citet{anttilaProsodic2009, anttilaMetrically2019} propose that the vowel alternations in the -r(ɪ) plurals are due to epenthetic vowels inserted to meet a word minimality requirement. I argue that this explanation does not hold in Northern Dagaare. Rather I propose that the singular suffix within the -r(ɪ) plural class is an empty mora. Additionally, I unify the -bɛ and -mɛ plural classes in Northern Dagaare and provide a preliminary analysis to explain the compensatory lengthening that occurs in the former but not the latter.
}


\IfFileExists{../localcommands.tex}{
   \addbibresource{../localbibliography.bib}
   %%%%%%%%%%%%%%%%%%%%%%%%%%%%%%%%%%%%%%%%%%%%%
%%%%%%%%%% From LSP %%%%%%%%%%%%%%%%%%%%%%%%%
%%%%%%%%%%%%%%%%%%%%%%%%%%%%%%%%%%%%%%%%%%%%%

\usepackage{tabularx,multicol}
\usepackage{url}
\urlstyle{same}

\usepackage{langsci-optional}
\usepackage{langsci-lgr}
\usepackage{langsci-gb4e}

% NOTE THAT THE FOLLOWING PACKAGES ARE BANNED: 
% - pstricks
% - tipa
%%%%%%%%%%%%%%%%%%%%%%%%%%%%%%%%%%%%%%%%%%%%%
%%%%%%%%%% From ACAL LaTeX template %%%%%%%%%
%%%%%%%%%%%%%%%%%%%%%%%%%%%%%%%%%%%%%%%%%%%%%

%For stacking text, used here in autosegmental diagrams
\usepackage{stackengine}

%To combine rows in tables
\usepackage{multirow}

%Gives some extra formatting options, e.g. underlining/strikeout
\usepackage[normalem]{ulem}

%Use package/command below to create a double-spaced document, if you want one. Uncomment BOTH the package and the command (\doublespacing) to create a doublespaced document, or leave them as is to have a single-spaced document.
%\usepackage{setspace}
%\doublespacing 

 

%use for special OT tableaux symbols like bomb and sad face. must be loaded early on because it doesn't play well with some other packages
\usepackage{fourier-orns}

%Basic math symbols 
\usepackage{pifont}
\usepackage{amssymb}

%%%Gives shortcuts for glossing. The use of this package is NOT explained in the Quick Reference Guide, but the documentation is on CTAN for those that are interested. MJKD finds it handy for glossing. (https://ctan.org/pkg/leipzig?lang=en)
\usepackage{leipzig}

%Tables
% \usepackage{caption} %For table captions 

%For OT-style tableaux
\usepackage[notipa]{ot-tableau}

%highlights text with \hl{text}
\usepackage{color, soul} %note that soul sometimes eats Unicode text witouth telling you. 

%Drawing Syntax Trees
\usepackage[linguistics]{forest}

%This specifies some formatting for the forest trees to make them nicer to look at. Unfortunately this was borrowed by Michael Diercks from someone a while ago and he can't remember who. Apologies and if someone lets him know he will update this.
\forestset{
  nice nodes/.style={
    for tree={
      inner sep=0pt,
      fit=band,
    },
  },
  default preamble=nice nodes,
}

%A couple of packages that seemed to prefer being called toward the end of the preamble

%This package provides macros for typesetting SPE-style phonological rules. I've redefined the \oneof command here, so that there the rule includes a right brace. 
\usepackage{phonrule}
\renewcommand{\oneof}[2][c]{\ensuremath{
    ~\left\{
    \begin{tabular}{#1#1}#2\end{tabular}
    \right\}
    }}

%For using Greek letters outside of math mode.
% \usepackage{textgreek}


%Random, lets us use the XeLaTeX logo. Not important to the template at all.
% \usepackage{metalogo}



%%%%%%%%%%%%%%%%%%%%%%%%%%%%%%%%%%%%%%%%%%%%%
%%%%%%%%%% From Smith paper %%%%%%%%%%%%%%%%%
%%%%%%%%%%%%%%%%%%%%%%%%%%%%%%%%%%%%%%%%%%%%%

\usetikzlibrary{positioning}
\usetikzlibrary{trees}
% \usepackage{pstricks} %pstricks is banned. Use tikz instead. 
% \usepackage{pst-node}
%\usepackage{tikz}

%%%%%%%%%%%%%%%%%%%%%%%%%%%%%%%%%%%%%%%%%%%%%
%%%%%%%%%% From Gluckman paper %%%%%%%%%%%%%%
%%%%%%%%%%%%%%%%%%%%%%%%%%%%%%%%%%%%%%%%%%%%%

\usepackage{amsmath}



%%%%%%%%%%%%%%%%%%%%%%%%%%%%%%%%%%%%%%%%%%%%%
%%%%%%%%%% From Kandybowicz paper %%%%%%%%%%%
%%%%%%%%%%%%%%%%%%%%%%%%%%%%%%%%%%%%%%%%%%%%%

%These are in the .tex, but there's nothing in the "Figures" folder so I'm not sure that it's actually doing anything. 
 
% \graphicspath{ {./Figures/} }

%%%%%%%%%%%%%%%%%%%%%%%%%%%%%%%%%%%%%%%%%%%%%
%%%%%%%%%% From Matte paper %%%%%%%%%%%%%%%%%
%%%%%%%%%%%%%%%%%%%%%%%%%%%%%%%%%%%%%%%%%%%%%

%%%%%%%%%%%%%%%%%%%%%%%%%%%%%%%%%%%%%%%%%%%%%
%%%%%%%%%% From Grollemund paper %%%%%%%%%%%%%%
%%%%%%%%%%%%%%%%%%%%%%%%%%%%%%%%%%%%%%%%%%%%%

% \usepackage{lscape}

%%%%%%%%%%%%%%%%%%%%%%%%%%%%%%%%%%%%%%%%%%%%%
%%%%%%%%%% From Burns paper %%%%%%%%%%%%%%
%%%%%%%%%%%%%%%%%%%%%%%%%%%%%%%%%%%%%%%%%%%%%


% \usepackage{url}
% \urlstyle{same}

\usepackage{listings}
\lstset{basicstyle=\ttfamily,tabsize=2,breaklines=true}

%MJKD commented out - these are standard for chapters in edited volumes, I don't think we need them here in the preamble. 

% \usepackage{tabularx,multicol}
% \usepackage{langsci-basic}
% \usepackage{langsci-gb4e}
% \bibliography{localbibliography}
% \newcommand{\orcid}[1]{}
% \pagenumbering{roman}

% \IfFileExists{../localcommands.tex}{%hack to check whether this is being compiled as part of a collection or standalone
%    %%%%%%%%%%%%%%%%%%%%%%%%%%%%%%%%%%%%%%%%%%%%%
%%%%%%%%%% From LSP %%%%%%%%%%%%%%%%%%%%%%%%%
%%%%%%%%%%%%%%%%%%%%%%%%%%%%%%%%%%%%%%%%%%%%%

\usepackage{tabularx,multicol}
\usepackage{url}
\urlstyle{same}

\usepackage{langsci-optional}
\usepackage{langsci-lgr}
\usepackage{langsci-gb4e}

% NOTE THAT THE FOLLOWING PACKAGES ARE BANNED: 
% - pstricks
% - tipa
%%%%%%%%%%%%%%%%%%%%%%%%%%%%%%%%%%%%%%%%%%%%%
%%%%%%%%%% From ACAL LaTeX template %%%%%%%%%
%%%%%%%%%%%%%%%%%%%%%%%%%%%%%%%%%%%%%%%%%%%%%

%For stacking text, used here in autosegmental diagrams
\usepackage{stackengine}

%To combine rows in tables
\usepackage{multirow}

%Gives some extra formatting options, e.g. underlining/strikeout
\usepackage[normalem]{ulem}

%Use package/command below to create a double-spaced document, if you want one. Uncomment BOTH the package and the command (\doublespacing) to create a doublespaced document, or leave them as is to have a single-spaced document.
%\usepackage{setspace}
%\doublespacing 

 

%use for special OT tableaux symbols like bomb and sad face. must be loaded early on because it doesn't play well with some other packages
\usepackage{fourier-orns}

%Basic math symbols 
\usepackage{pifont}
\usepackage{amssymb}

%%%Gives shortcuts for glossing. The use of this package is NOT explained in the Quick Reference Guide, but the documentation is on CTAN for those that are interested. MJKD finds it handy for glossing. (https://ctan.org/pkg/leipzig?lang=en)
\usepackage{leipzig}

%Tables
% \usepackage{caption} %For table captions 

%For OT-style tableaux
\usepackage[notipa]{ot-tableau}

%highlights text with \hl{text}
\usepackage{color, soul} %note that soul sometimes eats Unicode text witouth telling you. 

%Drawing Syntax Trees
\usepackage[linguistics]{forest}

%This specifies some formatting for the forest trees to make them nicer to look at. Unfortunately this was borrowed by Michael Diercks from someone a while ago and he can't remember who. Apologies and if someone lets him know he will update this.
\forestset{
  nice nodes/.style={
    for tree={
      inner sep=0pt,
      fit=band,
    },
  },
  default preamble=nice nodes,
}

%A couple of packages that seemed to prefer being called toward the end of the preamble

%This package provides macros for typesetting SPE-style phonological rules. I've redefined the \oneof command here, so that there the rule includes a right brace. 
\usepackage{phonrule}
\renewcommand{\oneof}[2][c]{\ensuremath{
    ~\left\{
    \begin{tabular}{#1#1}#2\end{tabular}
    \right\}
    }}

%For using Greek letters outside of math mode.
% \usepackage{textgreek}


%Random, lets us use the XeLaTeX logo. Not important to the template at all.
% \usepackage{metalogo}



%%%%%%%%%%%%%%%%%%%%%%%%%%%%%%%%%%%%%%%%%%%%%
%%%%%%%%%% From Smith paper %%%%%%%%%%%%%%%%%
%%%%%%%%%%%%%%%%%%%%%%%%%%%%%%%%%%%%%%%%%%%%%

\usetikzlibrary{positioning}
\usetikzlibrary{trees}
% \usepackage{pstricks} %pstricks is banned. Use tikz instead. 
% \usepackage{pst-node}
%\usepackage{tikz}

%%%%%%%%%%%%%%%%%%%%%%%%%%%%%%%%%%%%%%%%%%%%%
%%%%%%%%%% From Gluckman paper %%%%%%%%%%%%%%
%%%%%%%%%%%%%%%%%%%%%%%%%%%%%%%%%%%%%%%%%%%%%

\usepackage{amsmath}



%%%%%%%%%%%%%%%%%%%%%%%%%%%%%%%%%%%%%%%%%%%%%
%%%%%%%%%% From Kandybowicz paper %%%%%%%%%%%
%%%%%%%%%%%%%%%%%%%%%%%%%%%%%%%%%%%%%%%%%%%%%

%These are in the .tex, but there's nothing in the "Figures" folder so I'm not sure that it's actually doing anything. 
 
% \graphicspath{ {./Figures/} }

%%%%%%%%%%%%%%%%%%%%%%%%%%%%%%%%%%%%%%%%%%%%%
%%%%%%%%%% From Matte paper %%%%%%%%%%%%%%%%%
%%%%%%%%%%%%%%%%%%%%%%%%%%%%%%%%%%%%%%%%%%%%%

%%%%%%%%%%%%%%%%%%%%%%%%%%%%%%%%%%%%%%%%%%%%%
%%%%%%%%%% From Grollemund paper %%%%%%%%%%%%%%
%%%%%%%%%%%%%%%%%%%%%%%%%%%%%%%%%%%%%%%%%%%%%

% \usepackage{lscape}

%%%%%%%%%%%%%%%%%%%%%%%%%%%%%%%%%%%%%%%%%%%%%
%%%%%%%%%% From Burns paper %%%%%%%%%%%%%%
%%%%%%%%%%%%%%%%%%%%%%%%%%%%%%%%%%%%%%%%%%%%%


% \usepackage{url}
% \urlstyle{same}

\usepackage{listings}
\lstset{basicstyle=\ttfamily,tabsize=2,breaklines=true}

%MJKD commented out - these are standard for chapters in edited volumes, I don't think we need them here in the preamble. 

% \usepackage{tabularx,multicol}
% \usepackage{langsci-basic}
% \usepackage{langsci-gb4e}
% \bibliography{localbibliography}
% \newcommand{\orcid}[1]{}
% \pagenumbering{roman}

% \IfFileExists{../localcommands.tex}{%hack to check whether this is being compiled as part of a collection or standalone
%    %%%%%%%%%%%%%%%%%%%%%%%%%%%%%%%%%%%%%%%%%%%%%
%%%%%%%%%% From LSP %%%%%%%%%%%%%%%%%%%%%%%%%
%%%%%%%%%%%%%%%%%%%%%%%%%%%%%%%%%%%%%%%%%%%%%

\usepackage{tabularx,multicol}
\usepackage{url}
\urlstyle{same}

\usepackage{langsci-optional}
\usepackage{langsci-lgr}
\usepackage{langsci-gb4e}

% NOTE THAT THE FOLLOWING PACKAGES ARE BANNED: 
% - pstricks
% - tipa
%%%%%%%%%%%%%%%%%%%%%%%%%%%%%%%%%%%%%%%%%%%%%
%%%%%%%%%% From ACAL LaTeX template %%%%%%%%%
%%%%%%%%%%%%%%%%%%%%%%%%%%%%%%%%%%%%%%%%%%%%%

%For stacking text, used here in autosegmental diagrams
\usepackage{stackengine}

%To combine rows in tables
\usepackage{multirow}

%Gives some extra formatting options, e.g. underlining/strikeout
\usepackage[normalem]{ulem}

%Use package/command below to create a double-spaced document, if you want one. Uncomment BOTH the package and the command (\doublespacing) to create a doublespaced document, or leave them as is to have a single-spaced document.
%\usepackage{setspace}
%\doublespacing 

 

%use for special OT tableaux symbols like bomb and sad face. must be loaded early on because it doesn't play well with some other packages
\usepackage{fourier-orns}

%Basic math symbols 
\usepackage{pifont}
\usepackage{amssymb}

%%%Gives shortcuts for glossing. The use of this package is NOT explained in the Quick Reference Guide, but the documentation is on CTAN for those that are interested. MJKD finds it handy for glossing. (https://ctan.org/pkg/leipzig?lang=en)
\usepackage{leipzig}

%Tables
% \usepackage{caption} %For table captions 

%For OT-style tableaux
\usepackage[notipa]{ot-tableau}

%highlights text with \hl{text}
\usepackage{color, soul} %note that soul sometimes eats Unicode text witouth telling you. 

%Drawing Syntax Trees
\usepackage[linguistics]{forest}

%This specifies some formatting for the forest trees to make them nicer to look at. Unfortunately this was borrowed by Michael Diercks from someone a while ago and he can't remember who. Apologies and if someone lets him know he will update this.
\forestset{
  nice nodes/.style={
    for tree={
      inner sep=0pt,
      fit=band,
    },
  },
  default preamble=nice nodes,
}

%A couple of packages that seemed to prefer being called toward the end of the preamble

%This package provides macros for typesetting SPE-style phonological rules. I've redefined the \oneof command here, so that there the rule includes a right brace. 
\usepackage{phonrule}
\renewcommand{\oneof}[2][c]{\ensuremath{
    ~\left\{
    \begin{tabular}{#1#1}#2\end{tabular}
    \right\}
    }}

%For using Greek letters outside of math mode.
% \usepackage{textgreek}


%Random, lets us use the XeLaTeX logo. Not important to the template at all.
% \usepackage{metalogo}



%%%%%%%%%%%%%%%%%%%%%%%%%%%%%%%%%%%%%%%%%%%%%
%%%%%%%%%% From Smith paper %%%%%%%%%%%%%%%%%
%%%%%%%%%%%%%%%%%%%%%%%%%%%%%%%%%%%%%%%%%%%%%

\usetikzlibrary{positioning}
\usetikzlibrary{trees}
% \usepackage{pstricks} %pstricks is banned. Use tikz instead. 
% \usepackage{pst-node}
%\usepackage{tikz}

%%%%%%%%%%%%%%%%%%%%%%%%%%%%%%%%%%%%%%%%%%%%%
%%%%%%%%%% From Gluckman paper %%%%%%%%%%%%%%
%%%%%%%%%%%%%%%%%%%%%%%%%%%%%%%%%%%%%%%%%%%%%

\usepackage{amsmath}



%%%%%%%%%%%%%%%%%%%%%%%%%%%%%%%%%%%%%%%%%%%%%
%%%%%%%%%% From Kandybowicz paper %%%%%%%%%%%
%%%%%%%%%%%%%%%%%%%%%%%%%%%%%%%%%%%%%%%%%%%%%

%These are in the .tex, but there's nothing in the "Figures" folder so I'm not sure that it's actually doing anything. 
 
% \graphicspath{ {./Figures/} }

%%%%%%%%%%%%%%%%%%%%%%%%%%%%%%%%%%%%%%%%%%%%%
%%%%%%%%%% From Matte paper %%%%%%%%%%%%%%%%%
%%%%%%%%%%%%%%%%%%%%%%%%%%%%%%%%%%%%%%%%%%%%%

%%%%%%%%%%%%%%%%%%%%%%%%%%%%%%%%%%%%%%%%%%%%%
%%%%%%%%%% From Grollemund paper %%%%%%%%%%%%%%
%%%%%%%%%%%%%%%%%%%%%%%%%%%%%%%%%%%%%%%%%%%%%

% \usepackage{lscape}

%%%%%%%%%%%%%%%%%%%%%%%%%%%%%%%%%%%%%%%%%%%%%
%%%%%%%%%% From Burns paper %%%%%%%%%%%%%%
%%%%%%%%%%%%%%%%%%%%%%%%%%%%%%%%%%%%%%%%%%%%%


% \usepackage{url}
% \urlstyle{same}

\usepackage{listings}
\lstset{basicstyle=\ttfamily,tabsize=2,breaklines=true}

%MJKD commented out - these are standard for chapters in edited volumes, I don't think we need them here in the preamble. 

% \usepackage{tabularx,multicol}
% \usepackage{langsci-basic}
% \usepackage{langsci-gb4e}
% \bibliography{localbibliography}
% \newcommand{\orcid}[1]{}
% \pagenumbering{roman}

% \IfFileExists{../localcommands.tex}{%hack to check whether this is being compiled as part of a collection or standalone
%    \input{../localpackages}
%    \input{../localcommands}
%    \input{../localhyphenation}
%     \bibliography{localbibliography}
%     \togglepaper[23]
% }{}

% %custom footer for preprints
% \papernote{\scriptsize\normalfont
%     Change author.
%     Change title.
%     To appear in:
%     Change Volume Editor.  
%     Change volume title.  
%     Berlin: Language Science Press. [preliminary page numbering]
% }



%%%%%%%%%%%%%%%%%%%%%%%%%%%%%%%%%%%%%%%%%%%%%
%%%%%%%%%% From Feibig paper %%%%%%%%%%%%%%
%%%%%%%%%%%%%%%%%%%%%%%%%%%%%%%%%%%%%%%%%%%%%

\usepackage{tikz-qtree}
\usepackage{tikz-qtree-compat}

%%%%%%%%%%%%%%%%%%%%%%%%%%%%%%%%%%%%%%%%%%%%%
%%%%%%%%%% From Olson paper %%%%%%%%%%%%%%
%%%%%%%%%%%%%%%%%%%%%%%%%%%%%%%%%%%%%%%%%%%%%

%MJKD - TIPA is forbidden lol. Commenting out for now but we'll have to address this in the Olson paper

%\usepackage[tone]{tipa}

%%%%%%%%%%%%%%%%%%%%%%%%%%%%%%%%%%%%%%%%%%%%%
%%%%%%%%%% From Caragine paper %%%%%%%%%%%%%%
%%%%%%%%%%%%%%%%%%%%%%%%%%%%%%%%%%%%%%%%%%%%%

\usepackage{placeins}
% \usepackage{graphicx}
% \usepackage{float} %Overleaf suggested this as a solution


%%%%%%%%%%%%%%%%%%%%%%%%%%%%%%%%%%%%%%%%%%%%%
%%%%%%%%%% From Kieffer paper %%%%%%%%%%%%%%
%%%%%%%%%%%%%%%%%%%%%%%%%%%%%%%%%%%%%%%%%%%%%

\usepackage{array}

%%%%%%%%%%%%%%%%%%%%%%%%%%%%%%%%%%%%%%%%%%%%%
%%%%%%%%%% From Payne paper %%%%%%%%%%%%%%
%%%%%%%%%%%%%%%%%%%%%%%%%%%%%%%%%%%%%%%%%%%%%

\usepackage{enumerate}
\usepackage{array,ragged2e}
\usepackage{pgfplots}


%%%%%%%%%%%%%%%%%%%%%%%%%%%%%%%%%%%%%%%%%%%%%
%%%%%%%%%% From Diercks paper %%%%%%%%%%%%%%
%%%%%%%%%%%%%%%%%%%%%%%%%%%%%%%%%%%%%%%%%%%%%

% \usepackage{cgloss}



%%%%%%%%%%%%%%%%%%%%%%%%%%%%%%%%%%%%%%%%%%%%%
%%%%%%%%%% From Li paper %%%%%%%%%%%%%%
%%%%%%%%%%%%%%%%%%%%%%%%%%%%%%%%%%%%%%%%%%%%%



%%%%%%%%%%%%%%%%%%%%%%%%%%%%%%%%%%%%%%%%%%%%%
%%%%%%%%%% From Myers paper %%%%%%%%%%%%%%
%%%%%%%%%%%%%%%%%%%%%%%%%%%%%%%%%%%%%%%%%%%%%



\usepackage{tikz-qtree}
\usepackage{tikz-qtree-compat}


%Package that makes glossing slightly easier.
\RequirePackage{leipzig}
%\RequirePackage{mathtools} % for \mathrlap

% % % Shortcuts (various borrowed from Jason Zentz)
\RequirePackage{xspace}
\xspaceaddexceptions{]\}}
\usepackage{langsci-textipa}
\usepackage{langsci-branding}
\usepackage{tabto}

%    %%%%%%%%%%%%%%%%%%%%%%%%%%%%%%%%%%%%%%%%%%%%%
%%%%%%%%%% From LSP %%%%%%%%%%%%%%%%%%%%%%%%%
%%%%%%%%%%%%%%%%%%%%%%%%%%%%%%%%%%%%%%%%%%%%%


% \newcommand*{\orcid}{}


\makeatletter
\let\thetitle\@title
\let\theauthor\@author
\makeatother

\newcommand{\togglepaper}[1][0]{
   \bibliography{../localbibliography}
   \papernote{\scriptsize\normalfont
     \theauthor.
     \titleTemp.
     To appear in:
     E. Di Tor \& Herr Rausgeberin (ed.).
     Booktitle in localcommands.tex.
     Berlin: Language Science Press. [preliminary page numbering]
   }
   \pagenumbering{roman}
   \setcounter{chapter}{#1}
   \addtocounter{chapter}{-1}
}

\newbool{bookcompile}
\booltrue{bookcompile}
\newcommand{\bookorchapter}[2]{\ifbool{bookcompile}{#1}{#2}}


%%%%%%%%%%%%%%%%%%%%%%%%%%%%%%%%%%%%%%%%%%%%%
%%%%%%%%%% From ACAL LaTeX template %%%%%%%%%
%%%%%%%%%%%%%%%%%%%%%%%%%%%%%%%%%%%%%%%%%%%%%


\usepackage{tikz}

\newcommand{\circled}[1]{\begin{tikzpicture}[baseline=(word.base)]
\node[draw, rounded corners, text height=8pt, text depth=2pt, inner sep=2pt, outer sep=0pt, use as bounding box] (word) {#1};
\end{tikzpicture}
}


%%%%%%%%%%%%%%%%%%%%%%%%%%%%%%%%%%%%%%%%%%%%%%
%%%%%%%%%%%%%%%%%%%%%%%%%%%%%%%%%%%%%%%%%%%%%%

% Useful Ling Shortcuts


%This makes the \emptyset command be a nicer one
\let\oldemptyset\emptyset
\let\emptyset\varnothing
\newcommand{\nothing}{$\emptyset$}

%Not all of these are explained in the Quick Reference Guide, but they are here bc they are relevant to some of our students. 
\newcommand{\xbar}{\ensuremath{'}\xspace}
\newcommand{\xhead}{\ensuremath{^\circ}\xspace}
\newcommand{\Lb}[1]{$\text{[}_{\text{#1}}$ } %A more convenient left bracket
\newcommand{\Rb}[1]{$\text{]}_{\text{#1}}$ } %A more convenient left bracket
\newcommand{\gap}{\underline{\hspace{1.2em}}}
\newcommand{\vP}{\emph{v}P}
\newcommand{\lilv}{\emph{v}}
\newcommand{\Abar}{A$'$-} %A more convenient A-bar notation
\newcommand{\ph}{$\varphi$\xspace} %A more convenient phi
\newcommand{\pro}{\emph{pro}\xspace}
\newcommand{\subs}[1]{\textsubscript{#1}} %A more convenient subscript
\newcommand{\spells}{$\Longleftrightarrow$} %spellout arrow for morph spellout rules
\newcommand{\tr}[1]{\textit{t}\textsubscript{\textit{#1}}} %easy traces with subscript
\newcommand{\supers}[1]{\textsuperscript{#1}}

%These are borrowed/modified from Andrew Mackenzie's ling-macros package. We have copied them in here (instead of just using the package itself) to avoid a minor clash that was generating an error.


% Abbreviations for glossing, based on Leipzig
\newleipzig{hab}{hab}{habitual}
\newleipzig{rem}{rem}{remote}
\newleipzig{sm}{sm}{subject marker}
\newleipzig{t}{t}{tense}
\newleipzig{aa}{aa}{anti-agreement}
\newleipzig{pron}{pron}{pronoun}
\newleipzig{rec}{rec}{recent}
\newleipzig{om}{om}{object marker}
%\newleipzig{ipfv}{ipfv}{imperfective}
\newleipzig{asp}{asp}{aspect}
\newleipzig{lk}{lk}{linker}
\newleipzig{pcl}{pcl}{particle}
\newleipzig{stat}{stat}{stative}
\newleipzig{ints}{ints}{intensive}
\newleipzig{ascl}{ascl}{assertive subject clitic}
\newleipzig{nascl}{nascl}{non-assertive subject clitic}
\newleipzig{ta}{ta}{tense and/or aspect}
\newleipzig{assoc}{assoc}{associative marker}
\newleipzig{hon}{hon}{honorific}
%\newleipzig{whprt}{wh}{\wh particle}
\newleipzig{sa}{sa}{subject agreement}
\newleipzig{conj}{conj}{conjunction}
%\newleipzig{loc}{loc}{locative}
\newleipzig{expl}{expl}{expletive}
\newleipzig{rcm}{rcm}{reciprocal marker}
\newleipzig{pers}{pers}{persistive}
\newleipzig{fv}{fv}{final vowel}
%\newleipzig{}{}{} %this is just to copy for when I want to add more


%%%%%%%%%%%%%%%%%%%%%%%%%%%%%%%%%%%%%%%%%%%%%%%%%%%%%
%%%%%%%%%%%%%%%%%%%%%%%%%%%%%%%%%%%%%%%%%%%%%%%%%%%%%


%%%%%%%%%%%%%%%%%%%%%%%%%%%%%%%%%%%%%%%%%%%%%
%%%%%%%%%% From Gluckman paper %%%%%%%%%%%%%%
%%%%%%%%%%%%%%%%%%%%%%%%%%%%%%%%%%%%%%%%%%%%%

\newcommand{\pcite}[1]{\citeauthor{#1}'s (\citeyear{#1})}


%%%%%%%%%%%%%%%%%%%%%%%%%%%%%%%%%%%%%%%%%%%%%
%%%%%%%%%% From Myers paper %%%%%%%%%%%%%%
%%%%%%%%%%%%%%%%%%%%%%%%%%%%%%%%%%%%%%%%%%%%%

%% hspace over width of something without showing it
\newcommand*{\hspaceThis}[1]{\settowidth{\LSPTmp}{#1}\hspace*{\LSPTmp}}
% use \hphantom{} for this


%%%%%%%%%%%%%%%%%%%%%%%%%%%%%%%%%%%%%%%%%%%%%
%%%%%%%%%% From Matte paper %%%%%%%%%%%%%%%%%
%%%%%%%%%%%%%%%%%%%%%%%%%%%%%%%%%%%%%%%%%%%%%

%mjkd commented this out in case it's breaking something right now. 
% \newcounter{xlistex}
% \newcommand{\footnoteex}{\stepcounter{xlistex}(\roman{xlistex})}

\newleipzig{sp}{sp}{speaker} %a new leipzig glossing command


%%%%%%%%%%%%%%%%%%%%%%%%%%%%%%%%%%%%%%%%%%%%%
%%%%%%%%%% From GamFranich paper %%%%%%%%%%%%%%
%%%%%%%%%%%%%%%%%%%%%%%%%%%%%%%%%%%%%%%%%%%%%

%MJKD commented these out - both will be provided by the overall LSP book template 

% \bibliography{localbibliography}
% \newcommand{\orcid}[1]{}

%%%%%%%%%%%%%%%%%%%%%%%%%%%%%%%%%%%%%%%%%%%%%
%%%%%%%%%% From Diercks paper %%%%%%%%%%%%%%
%%%%%%%%%%%%%%%%%%%%%%%%%%%%%%%%%%%%%%%%%%%%%

\newcommand{\abar}{$\bar{\text{A}}$\xspace} % this one is different to the \Abar command in Mike's shortcuts doc
\newcommand{\topFeat}{[\textsc{top}]\xspace}


%%%%%%%%%%%%%%%%%%%%%%%%%%%%%%%%%%%%%%%%%%%%%
%%%%%%%%%% From Kinchsular paper %%%%%%%%%%%%%%
%%%%%%%%%%%%%%%%%%%%%%%%%%%%%%%%%%%%%%%%%%%%%

%%%%%%%%%%%%%%%%%%%%%%%%%%%%%%%%%%%%%%%%%%%%%
%%%%%%%%%% From Chen paper %%%%%%%%%%%%%%
%%%%%%%%%%%%%%%%%%%%%%%%%%%%%%%%%%%%%%%%%%%%%

\newcounter{savedex}
\makeatletter
\newcommand*{\saveex}{\setcounter{savedex}{\value{\@xnumctr}}}
\newcommand*{\resumeex}{\setcounter{\@xnumctr}{\value{savedex}}}
\makeatother


%    \input{../localhyphenation}
%     \bibliography{localbibliography}
%     \togglepaper[23]
% }{}

% %custom footer for preprints
% \papernote{\scriptsize\normalfont
%     Change author.
%     Change title.
%     To appear in:
%     Change Volume Editor.  
%     Change volume title.  
%     Berlin: Language Science Press. [preliminary page numbering]
% }



%%%%%%%%%%%%%%%%%%%%%%%%%%%%%%%%%%%%%%%%%%%%%
%%%%%%%%%% From Feibig paper %%%%%%%%%%%%%%
%%%%%%%%%%%%%%%%%%%%%%%%%%%%%%%%%%%%%%%%%%%%%

\usepackage{tikz-qtree}
\usepackage{tikz-qtree-compat}

%%%%%%%%%%%%%%%%%%%%%%%%%%%%%%%%%%%%%%%%%%%%%
%%%%%%%%%% From Olson paper %%%%%%%%%%%%%%
%%%%%%%%%%%%%%%%%%%%%%%%%%%%%%%%%%%%%%%%%%%%%

%MJKD - TIPA is forbidden lol. Commenting out for now but we'll have to address this in the Olson paper

%\usepackage[tone]{tipa}

%%%%%%%%%%%%%%%%%%%%%%%%%%%%%%%%%%%%%%%%%%%%%
%%%%%%%%%% From Caragine paper %%%%%%%%%%%%%%
%%%%%%%%%%%%%%%%%%%%%%%%%%%%%%%%%%%%%%%%%%%%%

\usepackage{placeins}
% \usepackage{graphicx}
% \usepackage{float} %Overleaf suggested this as a solution


%%%%%%%%%%%%%%%%%%%%%%%%%%%%%%%%%%%%%%%%%%%%%
%%%%%%%%%% From Kieffer paper %%%%%%%%%%%%%%
%%%%%%%%%%%%%%%%%%%%%%%%%%%%%%%%%%%%%%%%%%%%%

\usepackage{array}

%%%%%%%%%%%%%%%%%%%%%%%%%%%%%%%%%%%%%%%%%%%%%
%%%%%%%%%% From Payne paper %%%%%%%%%%%%%%
%%%%%%%%%%%%%%%%%%%%%%%%%%%%%%%%%%%%%%%%%%%%%

\usepackage{enumerate}
\usepackage{array,ragged2e}
\usepackage{pgfplots}


%%%%%%%%%%%%%%%%%%%%%%%%%%%%%%%%%%%%%%%%%%%%%
%%%%%%%%%% From Diercks paper %%%%%%%%%%%%%%
%%%%%%%%%%%%%%%%%%%%%%%%%%%%%%%%%%%%%%%%%%%%%

% \usepackage{cgloss}



%%%%%%%%%%%%%%%%%%%%%%%%%%%%%%%%%%%%%%%%%%%%%
%%%%%%%%%% From Li paper %%%%%%%%%%%%%%
%%%%%%%%%%%%%%%%%%%%%%%%%%%%%%%%%%%%%%%%%%%%%



%%%%%%%%%%%%%%%%%%%%%%%%%%%%%%%%%%%%%%%%%%%%%
%%%%%%%%%% From Myers paper %%%%%%%%%%%%%%
%%%%%%%%%%%%%%%%%%%%%%%%%%%%%%%%%%%%%%%%%%%%%



\usepackage{tikz-qtree}
\usepackage{tikz-qtree-compat}


%Package that makes glossing slightly easier.
\RequirePackage{leipzig}
%\RequirePackage{mathtools} % for \mathrlap

% % % Shortcuts (various borrowed from Jason Zentz)
\RequirePackage{xspace}
\xspaceaddexceptions{]\}}
\usepackage{langsci-textipa}
\usepackage{langsci-branding}
\usepackage{tabto}

%    %%%%%%%%%%%%%%%%%%%%%%%%%%%%%%%%%%%%%%%%%%%%%
%%%%%%%%%% From LSP %%%%%%%%%%%%%%%%%%%%%%%%%
%%%%%%%%%%%%%%%%%%%%%%%%%%%%%%%%%%%%%%%%%%%%%


% \newcommand*{\orcid}{}


\makeatletter
\let\thetitle\@title
\let\theauthor\@author
\makeatother

\newcommand{\togglepaper}[1][0]{
   \bibliography{../localbibliography}
   \papernote{\scriptsize\normalfont
     \theauthor.
     \titleTemp.
     To appear in:
     E. Di Tor \& Herr Rausgeberin (ed.).
     Booktitle in localcommands.tex.
     Berlin: Language Science Press. [preliminary page numbering]
   }
   \pagenumbering{roman}
   \setcounter{chapter}{#1}
   \addtocounter{chapter}{-1}
}

\newbool{bookcompile}
\booltrue{bookcompile}
\newcommand{\bookorchapter}[2]{\ifbool{bookcompile}{#1}{#2}}


%%%%%%%%%%%%%%%%%%%%%%%%%%%%%%%%%%%%%%%%%%%%%
%%%%%%%%%% From ACAL LaTeX template %%%%%%%%%
%%%%%%%%%%%%%%%%%%%%%%%%%%%%%%%%%%%%%%%%%%%%%


\usepackage{tikz}

\newcommand{\circled}[1]{\begin{tikzpicture}[baseline=(word.base)]
\node[draw, rounded corners, text height=8pt, text depth=2pt, inner sep=2pt, outer sep=0pt, use as bounding box] (word) {#1};
\end{tikzpicture}
}


%%%%%%%%%%%%%%%%%%%%%%%%%%%%%%%%%%%%%%%%%%%%%%
%%%%%%%%%%%%%%%%%%%%%%%%%%%%%%%%%%%%%%%%%%%%%%

% Useful Ling Shortcuts


%This makes the \emptyset command be a nicer one
\let\oldemptyset\emptyset
\let\emptyset\varnothing
\newcommand{\nothing}{$\emptyset$}

%Not all of these are explained in the Quick Reference Guide, but they are here bc they are relevant to some of our students. 
\newcommand{\xbar}{\ensuremath{'}\xspace}
\newcommand{\xhead}{\ensuremath{^\circ}\xspace}
\newcommand{\Lb}[1]{$\text{[}_{\text{#1}}$ } %A more convenient left bracket
\newcommand{\Rb}[1]{$\text{]}_{\text{#1}}$ } %A more convenient left bracket
\newcommand{\gap}{\underline{\hspace{1.2em}}}
\newcommand{\vP}{\emph{v}P}
\newcommand{\lilv}{\emph{v}}
\newcommand{\Abar}{A$'$-} %A more convenient A-bar notation
\newcommand{\ph}{$\varphi$\xspace} %A more convenient phi
\newcommand{\pro}{\emph{pro}\xspace}
\newcommand{\subs}[1]{\textsubscript{#1}} %A more convenient subscript
\newcommand{\spells}{$\Longleftrightarrow$} %spellout arrow for morph spellout rules
\newcommand{\tr}[1]{\textit{t}\textsubscript{\textit{#1}}} %easy traces with subscript
\newcommand{\supers}[1]{\textsuperscript{#1}}

%These are borrowed/modified from Andrew Mackenzie's ling-macros package. We have copied them in here (instead of just using the package itself) to avoid a minor clash that was generating an error.


% Abbreviations for glossing, based on Leipzig
\newleipzig{hab}{hab}{habitual}
\newleipzig{rem}{rem}{remote}
\newleipzig{sm}{sm}{subject marker}
\newleipzig{t}{t}{tense}
\newleipzig{aa}{aa}{anti-agreement}
\newleipzig{pron}{pron}{pronoun}
\newleipzig{rec}{rec}{recent}
\newleipzig{om}{om}{object marker}
%\newleipzig{ipfv}{ipfv}{imperfective}
\newleipzig{asp}{asp}{aspect}
\newleipzig{lk}{lk}{linker}
\newleipzig{pcl}{pcl}{particle}
\newleipzig{stat}{stat}{stative}
\newleipzig{ints}{ints}{intensive}
\newleipzig{ascl}{ascl}{assertive subject clitic}
\newleipzig{nascl}{nascl}{non-assertive subject clitic}
\newleipzig{ta}{ta}{tense and/or aspect}
\newleipzig{assoc}{assoc}{associative marker}
\newleipzig{hon}{hon}{honorific}
%\newleipzig{whprt}{wh}{\wh particle}
\newleipzig{sa}{sa}{subject agreement}
\newleipzig{conj}{conj}{conjunction}
%\newleipzig{loc}{loc}{locative}
\newleipzig{expl}{expl}{expletive}
\newleipzig{rcm}{rcm}{reciprocal marker}
\newleipzig{pers}{pers}{persistive}
\newleipzig{fv}{fv}{final vowel}
%\newleipzig{}{}{} %this is just to copy for when I want to add more


%%%%%%%%%%%%%%%%%%%%%%%%%%%%%%%%%%%%%%%%%%%%%%%%%%%%%
%%%%%%%%%%%%%%%%%%%%%%%%%%%%%%%%%%%%%%%%%%%%%%%%%%%%%


%%%%%%%%%%%%%%%%%%%%%%%%%%%%%%%%%%%%%%%%%%%%%
%%%%%%%%%% From Gluckman paper %%%%%%%%%%%%%%
%%%%%%%%%%%%%%%%%%%%%%%%%%%%%%%%%%%%%%%%%%%%%

\newcommand{\pcite}[1]{\citeauthor{#1}'s (\citeyear{#1})}


%%%%%%%%%%%%%%%%%%%%%%%%%%%%%%%%%%%%%%%%%%%%%
%%%%%%%%%% From Myers paper %%%%%%%%%%%%%%
%%%%%%%%%%%%%%%%%%%%%%%%%%%%%%%%%%%%%%%%%%%%%

%% hspace over width of something without showing it
\newcommand*{\hspaceThis}[1]{\settowidth{\LSPTmp}{#1}\hspace*{\LSPTmp}}
% use \hphantom{} for this


%%%%%%%%%%%%%%%%%%%%%%%%%%%%%%%%%%%%%%%%%%%%%
%%%%%%%%%% From Matte paper %%%%%%%%%%%%%%%%%
%%%%%%%%%%%%%%%%%%%%%%%%%%%%%%%%%%%%%%%%%%%%%

%mjkd commented this out in case it's breaking something right now. 
% \newcounter{xlistex}
% \newcommand{\footnoteex}{\stepcounter{xlistex}(\roman{xlistex})}

\newleipzig{sp}{sp}{speaker} %a new leipzig glossing command


%%%%%%%%%%%%%%%%%%%%%%%%%%%%%%%%%%%%%%%%%%%%%
%%%%%%%%%% From GamFranich paper %%%%%%%%%%%%%%
%%%%%%%%%%%%%%%%%%%%%%%%%%%%%%%%%%%%%%%%%%%%%

%MJKD commented these out - both will be provided by the overall LSP book template 

% \bibliography{localbibliography}
% \newcommand{\orcid}[1]{}

%%%%%%%%%%%%%%%%%%%%%%%%%%%%%%%%%%%%%%%%%%%%%
%%%%%%%%%% From Diercks paper %%%%%%%%%%%%%%
%%%%%%%%%%%%%%%%%%%%%%%%%%%%%%%%%%%%%%%%%%%%%

\newcommand{\abar}{$\bar{\text{A}}$\xspace} % this one is different to the \Abar command in Mike's shortcuts doc
\newcommand{\topFeat}{[\textsc{top}]\xspace}


%%%%%%%%%%%%%%%%%%%%%%%%%%%%%%%%%%%%%%%%%%%%%
%%%%%%%%%% From Kinchsular paper %%%%%%%%%%%%%%
%%%%%%%%%%%%%%%%%%%%%%%%%%%%%%%%%%%%%%%%%%%%%

%%%%%%%%%%%%%%%%%%%%%%%%%%%%%%%%%%%%%%%%%%%%%
%%%%%%%%%% From Chen paper %%%%%%%%%%%%%%
%%%%%%%%%%%%%%%%%%%%%%%%%%%%%%%%%%%%%%%%%%%%%

\newcounter{savedex}
\makeatletter
\newcommand*{\saveex}{\setcounter{savedex}{\value{\@xnumctr}}}
\newcommand*{\resumeex}{\setcounter{\@xnumctr}{\value{savedex}}}
\makeatother


%    \input{../localhyphenation}
%     \bibliography{localbibliography}
%     \togglepaper[23]
% }{}

% %custom footer for preprints
% \papernote{\scriptsize\normalfont
%     Change author.
%     Change title.
%     To appear in:
%     Change Volume Editor.  
%     Change volume title.  
%     Berlin: Language Science Press. [preliminary page numbering]
% }



%%%%%%%%%%%%%%%%%%%%%%%%%%%%%%%%%%%%%%%%%%%%%
%%%%%%%%%% From Feibig paper %%%%%%%%%%%%%%
%%%%%%%%%%%%%%%%%%%%%%%%%%%%%%%%%%%%%%%%%%%%%

\usepackage{tikz-qtree}
\usepackage{tikz-qtree-compat}

%%%%%%%%%%%%%%%%%%%%%%%%%%%%%%%%%%%%%%%%%%%%%
%%%%%%%%%% From Olson paper %%%%%%%%%%%%%%
%%%%%%%%%%%%%%%%%%%%%%%%%%%%%%%%%%%%%%%%%%%%%

%MJKD - TIPA is forbidden lol. Commenting out for now but we'll have to address this in the Olson paper

%\usepackage[tone]{tipa}

%%%%%%%%%%%%%%%%%%%%%%%%%%%%%%%%%%%%%%%%%%%%%
%%%%%%%%%% From Caragine paper %%%%%%%%%%%%%%
%%%%%%%%%%%%%%%%%%%%%%%%%%%%%%%%%%%%%%%%%%%%%

\usepackage{placeins}
% \usepackage{graphicx}
% \usepackage{float} %Overleaf suggested this as a solution


%%%%%%%%%%%%%%%%%%%%%%%%%%%%%%%%%%%%%%%%%%%%%
%%%%%%%%%% From Kieffer paper %%%%%%%%%%%%%%
%%%%%%%%%%%%%%%%%%%%%%%%%%%%%%%%%%%%%%%%%%%%%

\usepackage{array}

%%%%%%%%%%%%%%%%%%%%%%%%%%%%%%%%%%%%%%%%%%%%%
%%%%%%%%%% From Payne paper %%%%%%%%%%%%%%
%%%%%%%%%%%%%%%%%%%%%%%%%%%%%%%%%%%%%%%%%%%%%

\usepackage{enumerate}
\usepackage{array,ragged2e}
\usepackage{pgfplots}


%%%%%%%%%%%%%%%%%%%%%%%%%%%%%%%%%%%%%%%%%%%%%
%%%%%%%%%% From Diercks paper %%%%%%%%%%%%%%
%%%%%%%%%%%%%%%%%%%%%%%%%%%%%%%%%%%%%%%%%%%%%

% \usepackage{cgloss}



%%%%%%%%%%%%%%%%%%%%%%%%%%%%%%%%%%%%%%%%%%%%%
%%%%%%%%%% From Li paper %%%%%%%%%%%%%%
%%%%%%%%%%%%%%%%%%%%%%%%%%%%%%%%%%%%%%%%%%%%%



%%%%%%%%%%%%%%%%%%%%%%%%%%%%%%%%%%%%%%%%%%%%%
%%%%%%%%%% From Myers paper %%%%%%%%%%%%%%
%%%%%%%%%%%%%%%%%%%%%%%%%%%%%%%%%%%%%%%%%%%%%



\usepackage{tikz-qtree}
\usepackage{tikz-qtree-compat}


%Package that makes glossing slightly easier.
\RequirePackage{leipzig}
%\RequirePackage{mathtools} % for \mathrlap

% % % Shortcuts (various borrowed from Jason Zentz)
\RequirePackage{xspace}
\xspaceaddexceptions{]\}}
\usepackage{langsci-textipa}
\usepackage{langsci-branding}
\usepackage{tabto}

   %%%%%%%%%%%%%%%%%%%%%%%%%%%%%%%%%%%%%%%%%%%%%
%%%%%%%%%% From LSP %%%%%%%%%%%%%%%%%%%%%%%%%
%%%%%%%%%%%%%%%%%%%%%%%%%%%%%%%%%%%%%%%%%%%%%


% \newcommand*{\orcid}{}


\makeatletter
\let\thetitle\@title
\let\theauthor\@author
\makeatother

\newcommand{\togglepaper}[1][0]{
   \bibliography{../localbibliography}
   \papernote{\scriptsize\normalfont
     \theauthor.
     \titleTemp.
     To appear in:
     E. Di Tor \& Herr Rausgeberin (ed.).
     Booktitle in localcommands.tex.
     Berlin: Language Science Press. [preliminary page numbering]
   }
   \pagenumbering{roman}
   \setcounter{chapter}{#1}
   \addtocounter{chapter}{-1}
}

\newbool{bookcompile}
\booltrue{bookcompile}
\newcommand{\bookorchapter}[2]{\ifbool{bookcompile}{#1}{#2}}


%%%%%%%%%%%%%%%%%%%%%%%%%%%%%%%%%%%%%%%%%%%%%
%%%%%%%%%% From ACAL LaTeX template %%%%%%%%%
%%%%%%%%%%%%%%%%%%%%%%%%%%%%%%%%%%%%%%%%%%%%%


\usepackage{tikz}

\newcommand{\circled}[1]{\begin{tikzpicture}[baseline=(word.base)]
\node[draw, rounded corners, text height=8pt, text depth=2pt, inner sep=2pt, outer sep=0pt, use as bounding box] (word) {#1};
\end{tikzpicture}
}


%%%%%%%%%%%%%%%%%%%%%%%%%%%%%%%%%%%%%%%%%%%%%%
%%%%%%%%%%%%%%%%%%%%%%%%%%%%%%%%%%%%%%%%%%%%%%

% Useful Ling Shortcuts


%This makes the \emptyset command be a nicer one
\let\oldemptyset\emptyset
\let\emptyset\varnothing
\newcommand{\nothing}{$\emptyset$}

%Not all of these are explained in the Quick Reference Guide, but they are here bc they are relevant to some of our students. 
\newcommand{\xbar}{\ensuremath{'}\xspace}
\newcommand{\xhead}{\ensuremath{^\circ}\xspace}
\newcommand{\Lb}[1]{$\text{[}_{\text{#1}}$ } %A more convenient left bracket
\newcommand{\Rb}[1]{$\text{]}_{\text{#1}}$ } %A more convenient left bracket
\newcommand{\gap}{\underline{\hspace{1.2em}}}
\newcommand{\vP}{\emph{v}P}
\newcommand{\lilv}{\emph{v}}
\newcommand{\Abar}{A$'$-} %A more convenient A-bar notation
\newcommand{\ph}{$\varphi$\xspace} %A more convenient phi
\newcommand{\pro}{\emph{pro}\xspace}
\newcommand{\subs}[1]{\textsubscript{#1}} %A more convenient subscript
\newcommand{\spells}{$\Longleftrightarrow$} %spellout arrow for morph spellout rules
\newcommand{\tr}[1]{\textit{t}\textsubscript{\textit{#1}}} %easy traces with subscript
\newcommand{\supers}[1]{\textsuperscript{#1}}

%These are borrowed/modified from Andrew Mackenzie's ling-macros package. We have copied them in here (instead of just using the package itself) to avoid a minor clash that was generating an error.


% Abbreviations for glossing, based on Leipzig
\newleipzig{hab}{hab}{habitual}
\newleipzig{rem}{rem}{remote}
\newleipzig{sm}{sm}{subject marker}
\newleipzig{t}{t}{tense}
\newleipzig{aa}{aa}{anti-agreement}
\newleipzig{pron}{pron}{pronoun}
\newleipzig{rec}{rec}{recent}
\newleipzig{om}{om}{object marker}
%\newleipzig{ipfv}{ipfv}{imperfective}
\newleipzig{asp}{asp}{aspect}
\newleipzig{lk}{lk}{linker}
\newleipzig{pcl}{pcl}{particle}
\newleipzig{stat}{stat}{stative}
\newleipzig{ints}{ints}{intensive}
\newleipzig{ascl}{ascl}{assertive subject clitic}
\newleipzig{nascl}{nascl}{non-assertive subject clitic}
\newleipzig{ta}{ta}{tense and/or aspect}
\newleipzig{assoc}{assoc}{associative marker}
\newleipzig{hon}{hon}{honorific}
%\newleipzig{whprt}{wh}{\wh particle}
\newleipzig{sa}{sa}{subject agreement}
\newleipzig{conj}{conj}{conjunction}
%\newleipzig{loc}{loc}{locative}
\newleipzig{expl}{expl}{expletive}
\newleipzig{rcm}{rcm}{reciprocal marker}
\newleipzig{pers}{pers}{persistive}
\newleipzig{fv}{fv}{final vowel}
%\newleipzig{}{}{} %this is just to copy for when I want to add more


%%%%%%%%%%%%%%%%%%%%%%%%%%%%%%%%%%%%%%%%%%%%%%%%%%%%%
%%%%%%%%%%%%%%%%%%%%%%%%%%%%%%%%%%%%%%%%%%%%%%%%%%%%%


%%%%%%%%%%%%%%%%%%%%%%%%%%%%%%%%%%%%%%%%%%%%%
%%%%%%%%%% From Gluckman paper %%%%%%%%%%%%%%
%%%%%%%%%%%%%%%%%%%%%%%%%%%%%%%%%%%%%%%%%%%%%

\newcommand{\pcite}[1]{\citeauthor{#1}'s (\citeyear{#1})}


%%%%%%%%%%%%%%%%%%%%%%%%%%%%%%%%%%%%%%%%%%%%%
%%%%%%%%%% From Myers paper %%%%%%%%%%%%%%
%%%%%%%%%%%%%%%%%%%%%%%%%%%%%%%%%%%%%%%%%%%%%

%% hspace over width of something without showing it
\newcommand*{\hspaceThis}[1]{\settowidth{\LSPTmp}{#1}\hspace*{\LSPTmp}}
% use \hphantom{} for this


%%%%%%%%%%%%%%%%%%%%%%%%%%%%%%%%%%%%%%%%%%%%%
%%%%%%%%%% From Matte paper %%%%%%%%%%%%%%%%%
%%%%%%%%%%%%%%%%%%%%%%%%%%%%%%%%%%%%%%%%%%%%%

%mjkd commented this out in case it's breaking something right now. 
% \newcounter{xlistex}
% \newcommand{\footnoteex}{\stepcounter{xlistex}(\roman{xlistex})}

\newleipzig{sp}{sp}{speaker} %a new leipzig glossing command


%%%%%%%%%%%%%%%%%%%%%%%%%%%%%%%%%%%%%%%%%%%%%
%%%%%%%%%% From GamFranich paper %%%%%%%%%%%%%%
%%%%%%%%%%%%%%%%%%%%%%%%%%%%%%%%%%%%%%%%%%%%%

%MJKD commented these out - both will be provided by the overall LSP book template 

% \bibliography{localbibliography}
% \newcommand{\orcid}[1]{}

%%%%%%%%%%%%%%%%%%%%%%%%%%%%%%%%%%%%%%%%%%%%%
%%%%%%%%%% From Diercks paper %%%%%%%%%%%%%%
%%%%%%%%%%%%%%%%%%%%%%%%%%%%%%%%%%%%%%%%%%%%%

\newcommand{\abar}{$\bar{\text{A}}$\xspace} % this one is different to the \Abar command in Mike's shortcuts doc
\newcommand{\topFeat}{[\textsc{top}]\xspace}


%%%%%%%%%%%%%%%%%%%%%%%%%%%%%%%%%%%%%%%%%%%%%
%%%%%%%%%% From Kinchsular paper %%%%%%%%%%%%%%
%%%%%%%%%%%%%%%%%%%%%%%%%%%%%%%%%%%%%%%%%%%%%

%%%%%%%%%%%%%%%%%%%%%%%%%%%%%%%%%%%%%%%%%%%%%
%%%%%%%%%% From Chen paper %%%%%%%%%%%%%%
%%%%%%%%%%%%%%%%%%%%%%%%%%%%%%%%%%%%%%%%%%%%%

\newcounter{savedex}
\makeatletter
\newcommand*{\saveex}{\setcounter{savedex}{\value{\@xnumctr}}}
\newcommand*{\resumeex}{\setcounter{\@xnumctr}{\value{savedex}}}
\makeatother


   \input{../localhyphenation}
   \boolfalse{bookcompile}
   \togglepaper[23]%%chapternumber
}{}

\begin{document}
\maketitle

\section{Introduction}
\label{sec:caragine:intro}

The Dagaare (Mabia)\footnote{Dagaare is often also called Dagaari, Dagara, Waale, and Birifor. The Mabia language family is sometimes referred to as Gur.} language is made up of a dialect continuum which stretches from Northern Ghana to southern Burkina Faso. In this paper, I present novel data from the Nandome dialect of Northern Dagaare, which exhibits vowel alternations between several singular and plural nouns. \citet{anttilaProsodic2009, anttilaMetrically2019} found that similar patterns in Central Dagaare are driven by metrical foot structure, minimal prosodic word requirements and other language general constraints. I investigate whether this analysis can be extended to data from Northern Dagaare. I argue that despite their success capturing the variation in Central Dagaare, this question in the larger Dagaare language is far from settled. While the Anttila and Bodomo analysis makes the correct predictions for the -r singulars, it does not accurately explain the variation seen in the -r(ɪ) plurals and it does not address vowel length alternations in the -bɛ plurals. 

I maintain a simple affixation account of the -r singulars put forth by \citet{anttilaProsodic2009}, but I propose that within the -r(ɪ) plural declension class, the singulars receive a suffix consisting only of an empty mora. This mora is then filled by linking it to the features of the preceding vowel, resulting in vowel lengthening. Then language general constraints adopted from \citet{anttilaProsodic2009} sometimes cause diphthongization. In the -bɛ plurals, I argue that there is compensatory lengthening of the stem vowel in the plural when the stem final consonant is deleted. When the final consonant is a nasal it coalesces with the [b] in the suffix to form an [m]. This [m] is ambisyllabic and thus no compensatory lengthening is triggered.

The paper is organized as follows. \sectref{sec:caragine:Dagaare} introduces the Northern Dagaare dialect and its relevant linguistic features,  \sectref{sec:caragine:numberMarking} presents data of number marking, \sectref{sec:caragine:analysis} works towards an analysis, and \sectref{sec:caragine:conclusion} concludes.

\section{Northern Dagaare} \label{sec:caragine:Dagaare}

\subsection{Language background} \label{sec:caragine:langback}

Dagaare is an understudied and under-documented Mabia language (Niger-Congo) spoken in northwestern Ghana and southern Burkina Faso by between 1.5 and 2 million speakers (\citealp{aliDictionary2021,lewisEthnologue2009}). 
The Dagaare language is made up of a dialect continuum that, roughly speaking, can be divided into four main dialect groups -- Northern, Southern, Western and Central -- which show substantial variation  (\citealp{bodomoStructure1997, angsongna2022, aliDictionary2021}). The map in \figref{fig:caragine:map} shows the locations in which the dialect groups are spoken. This paper focuses on Nandome Dagaare, an under-documented sub-dialect of Northern Dagaare spoken in and around Nandom, Ghana.

\begin{figure}[htp]
    \centering
    \includegraphics[width=10cm]{figures/Dagaare-map.png}
    \caption{Map of the area where Dagaare is spoken \citep{angsongna2022}}
    \label{fig:caragine:map}
\end{figure}

A team of researchers at Georgetown University\footnote{Michael Obiri-Yeboah, Sophie Migacz, and Jack Pruett} and the University of Maryland\footnote{Cassandra Caragine} has spent over a year working with a native speaker to document and analyze this understudied dialect. The speaker is a male (27), who grew up in Kumasi, Ghana and speaks the Nandome dialect of Dagaare. Additionally, he speaks Twi, English, French, and Spanish. All interview sessions were recorded (about 40 hours) with a Zoom H4n Pro handheld recorder. Data was transcribed and uploaded to Twisted Tongues database management software. Researchers used Praat \citep{praatCaragine} for confirmation and correction of initial transcriptions. This resulted in 68 singular-plural pairs in Northern Dagaare.

\subsection{Phonological background} \label{sec:Caragine:phonback}

Northern Dagaare (ND) has a nine-vowel system consisting of four +/-\textsc{atr} (Advanced Tongue Root)  pairs and one low central vowel, \tabref{tab:caragine:VowelChart}. The low central vowel /a/ is likely under-specified for \textsc{atr}; evidence for this is presented later in this section.  Vowels also contrast in length ([nà] ‘this’ v. [nàà] ‘chief’), tone ([báá] ‘dog’ v. [bàà] ‘river’), and nasality ([sàà] ‘rain’ v. [sã̀ã̀] ‘father’). 

\begin{table}
\begin{tabular}{ccccc}
    \lsptoprule 
  {}  & \multicolumn{2}{c}{Front} & \multicolumn{2}{c}{Back} \\
  {}  & +\textsc{atr} &-\textsc{atr} &+\textsc{atr} &-\textsc{atr} \\
    \midrule \midrule
    High & i & ɪ & u & ʊ \\
    Mid & e & ɛ & o & ɔ \\
    Low & \multicolumn{4}{c}{a} \\
    \lspbottomrule 
\end{tabular}
\caption{Dagaare Vowel Chart} \label{tab:caragine:VowelChart}
\end{table}

Additionally, ND exhibits \textsc{atr} vowel harmony. In the nominal domain, this process is progressive. The low central vowel /a/ is invisible to this process; it does not undergo, trigger, or block vowel harmony. Thus, it is likely that /a/ is underspecified for \textsc{atr}. For further discussion of vowel harmony in  Dagaare see \citet{bodomoStructure1997} and \citet{aliDictionary2021}.

 Northern Dagaare has a (C)V(V)(C) syllable structure. It allows optional onsets and optional codas as well as short vowels, long vowels, and diphthongs. Complex onsets are currently unattested, and there is only one word in our data set which contains a complex coda: [tʊ̀βr] 'ears'. Syllable weight and moraic structure will be considered in \sectref{sec:caragine:bepluralanalysis}.

%%%%%%%%%%%%%%%%%%%%%%%%%%%%%%%%%%%%%%%%%%%%%%%%%%%%%%%%%%%%%
\section{Number marking} \label{sec:caragine:numberMarking}

The Dagaare nominal system has two numbers: singular and plural. Dagaare is noted to have an `inverse' or `polarity' system, in which a single morpheme can act as either a singular or plural marker \citep{Grimm2012, Grimm2021, anttilaProsodic2009}. As \tabref{tab:Caragine:polarity} shows, -r can mark either singulars or plurals depending on the declension class of the noun.

\begin{table}
  \caption{Polarity Marking}
  \label{tab:Caragine:polarity}
\begin{tabular}{lll}
\lsptoprule 
Singular & Plural & Gloss\\
\midrule\midrule
nóbíè & nóbí-r & `finger'\\
jí-r & jíé & `house'\\
\lspbottomrule 
\end{tabular}
\end{table}

However, -r is not the only number marker. There are several ways that number can be expressed on the noun depending on its declension class. The declension classes of Dagaare have been categorized in several ways. I use the six classes shown in \tabref{tab:Caragine:classes}, but the exact categorization is tangential to the current paper. Additionally, the labels used for each class have not be standardized and are provided for referential simplicity only. As noted in \sectref{sec:Caragine:phonback}, vowel harmony is an active process in the language. The citation forms of each of the suffixes is given in the -\textsc{atr} form as this is what occurs with stems whose only vowel is /a/, which is underspecified for \textsc{atr}. The specific rules governing which number markers are used for a particular noun are outside the scope of this paper. For further discussion of this topic see \citet{Pruett}. 

\begin{table}
  \caption{ND Declension Classes}
  \label{tab:Caragine:classes}
\begin{tabular}{llll}
\lsptoprule 
& Singular & Plural & Gloss\\
\midrule\midrule
-r singulars & jí-r & jí-é & `house' \\
-r(ɪ) plurals &  bàà & bàà-r & `river' \\
-bɛ plurals & dɛ́b̚ & dɛ̀ɛ̀-βɛ̀ & `man'\\
l nominals &jéèpùlè &	jéépùlì & `sister'\\
-mɪnɛ plurals & sã̀ã̀ & sã̀ã̀-mɪ̀nɛ́ & `father'\\
irregulars & wáàb &	wɪ́ɪ́r &	`snake' \\

\lspbottomrule 
\end{tabular}
\end{table}

The declension classes and a representative example for each are included in \ref{tab:Caragine:classes}. The first three of these classes (-r singulars, -r(ɪ) plurals, and -bɛ plurals) exhibit some sort of vowel alternation. Sections \ref{sec:Caragine:ePlural}-\ref{sec:Caragine:bePl} will lay out the patterns seen in these noun classes in Northern Dagaare and compare them to their correspondents in Central Dagaare (CD).  

%%%%%%%%%%%%%%%%%%%%%%%%%%%%%
\subsection{-r singulars} \label{sec:Caragine:ePlural}

The -r singular declension class is one of the classes exhibiting polarity marking with -r(ɪ). As can be seen in \tabref{tab:Caragine:rSing}, this class overtly marks both the singular and plural with a suffix. An -r is suffixed in the singular and -ɛ is suffixed in the plural. The -ɛ plural suffix undergoes \textsc{atr} vowel harmony to agree with the stem. When the -ɛ suffix results in vowel hiatus, there is no rearticulation. When the qualities are different for the two vowels, the phonetic yield is a diphthong, and when the two vowels have the same quality (eg. when -ɛ is added to an [ɛ] or [e] final stem), a long vowel is produced.

The left 2 columns of \tabref{tab:Caragine:rSing} provide novel singular/plural pairs in the -r singular declension class of Northern Dagaare. The right two columns provide the corresponding word in Central Dagaare.\footnote{In the original sources standard Dagaare orthography was used instead of IPA. For consistency and ease of comparison, the orthography has been converted to IPA. There is generally a one-to-one correspondence between the orthography and IPA. The only exceptions to this are <e> which can be either [e] or [ɪ] and <o> which can be either [o] or [ʊ]. I used whichever IPA symbol agreed in \textsc{atr} with the rest of the word. When the context was ambiguous, I selected the symbol which more closely resembled the ND pronunciation. Tones are not recoverable from the orthography, so they have been omitted. These reconstructed examples have been marked with an asterisk (*). A dash (--) indicates that a translation could not be found for this word. When there is a semantic difference between the ND and CD words, the CD meaning is indicated in parenthesis.} 
Unless otherwise noted, here and throughout the paper, the Central Dagaare data come from \citet{aliDictionary2021}. 

\begin{table}

  \caption{-r singulars in Northern Dagaare and corresponding nouns in Central Dagaare. Central Dagaare data from \citet{aliDictionary2021}}
  \label{tab:Caragine:rSing}
  
\begin{tabular}{lllll}
\lsptoprule 
\multicolumn{2}{c}{Northern Dagaare} & \multicolumn{1}{c}{}& \multicolumn{2}{c}{Central Dagaare}\\
Singular & Plural & Gloss 
& Singular & Plural \\
\midrule\midrule
bàpɛ́-r & bàpɛ́-ɛ́ & `male friend' &bábí-rì& *babi-e (`best friend')\\
g͡bɛ́-r & g͡bɛ́-ɛ́ & `leg' & *g͡bɛ́-rɪ̀ & *g͡bɛ-ɛ\\
tʊ́β-r & tʊ̀β-ɛ́ & `ear' & tòò-rí & *tobo \\
jí-r & jí-é & `house' & jí-rì& *ji-e\\
dɛ̀-r & dɛ̀-ɛ̀ & `stair' &--&--\\
bũ̀bí-r & bũ̀bì-é & `seed' & bʊ́mbí-rì & *bombi-e\\  
mʊ̀ndɛ́-r & mʊ̀ndɛ́-ɛ́ & `spider' & bádɛ́-rɪ́ & *badɛr-rɪ \\
màrbí-r & màrbí-è & `star' & ŋ͡márɪ́bí-rì & *ŋ͡marebi-e \\
\lspbottomrule 

\end{tabular}
\end{table}

As can be seen in \tabref{tab:Caragine:rSing}, the the singular -r suffix in Northern Dagaare always corresponds to a -rɪ suffix in Central Dagaare. `Ear' and `spider' appear to be irregular plurals in CD, but the remaining CD plurals are marked with a corresponding suffix to that used in ND.

%%%%%%%%%%%%%%%%%%%%%%%%%%%%%%%%%%%%%%%%%
\subsection{-r(ɪ) plurals} \label{sec:caragine:rPlurals}
\largerpage
The most common of the declension classes adds -r or -rɪ in the plural (18/68 noun pairs). Exactly what conditions the final [ɪ] is still unknown.\footnote{There may also be some synchronic variation. When used in a sentence context rather than in isolation, those forms which still ``retain'' the final [ɪ] may reduce/delete the final vowel [tɪ́-ɾɪ́] \sim [tɪ́r] `tree'. This can alternate within a single speaker from one utterance to another. It is possible that since the -r is the only information needed to distinguish the singular from the plural,  Northern Dagaare may be losing this final vowel in more common nouns. Further research is required to determine exactly where and why the final [ɪ] is sometimes dropped.}
Other dialects seem to only have -rɪ \citep{anttilaProsodic2009, anttilaMetrically2019, bodomo2007} -- see Tables \ref{tab:Caragine:rPluralsSame}-\ref{tab:Caragine:rPluralsChange} below for examples in CD. Interestingly, as was seen in \sectref{sec:Caragine:ePlural}, the [ɪ] has universally been lost when it marks the singular of the -r singular declension class even though their correlate in other dialecs always has a final [ɪ]. \tabref{tab:Caragine:rPluralsSame} shows some examples of plurals which are formed by the addition of -r(ɪ).

\begin{table}

  \caption{-r(ɪ) plurals in Northern Dagaare and corresponding nouns in Central Dagaare} 
  \label{tab:Caragine:rPluralsSame}
  
\begin{tabular}{lllll}
\lsptoprule 
\multicolumn{2}{c}{Northern Dagaare} & \multicolumn{1}{c}{}& \multicolumn{2}{c}{Central Dagaare}\\
Singular & Plural & Gloss 
& Singular & Plural \\
\midrule\midrule
báá & báá-r & `dog' & báá & *báá-rí \citep{bodomoStructure1997}\\
bàà & bàà-r & `river' & bàá & *bà-rí \citep{bodomoStructure1997}\\
dáá & dáá-r & `stick' & dàá & *dà-rí \citep{bodomoStructure1997}\\
nàfáʊ̀ & nàfáʊ̀-ɾɪ́ & `shoe' & *nɔɔtɪ-rʊʊ & *nɔɔtɪ-ɛ \\
gàdò & gàdò-ɾí & `bed' & gádó & *gado-ri \\
mʊ̀ɔ́ & mʊ̀ɔ́-r& `grass' & mɔ́-rɪ́ & *mɔ-ɛ\\
\lspbottomrule 

\end{tabular}
\end{table}

As seen in \tabref{tab:Caragine:rPluralsSame}, in this declension class the singulars are unmarked and the plurals are marked by either -r or -ɾɪ́.\footnote{Note that in Tables \ref{tab:Caragine:rPluralsSame}-\ref{tab:Caragine:rPluralsChange} all stems which take the -r(ɪ) plural end in vowels. For discussion of how this may impact declension class, see \citet{Pruett}} When the plural marker has a final [ɪ], there is a general allophonic process in ND in which /r/ becomes [ɾ] in syllable onsets. 

As can be seen in \tabref{tab:Caragine:rPluralsChange}, some nouns in the -r(ɪ) plural declension class show variations in the vowels between the singular and plural forms. The singular nouns in ND have a vowel which does not appear in the plural. In some words, like [núú], the extra vowel in the singular is the same as the vowel in the plural stem [nú], but in others, for example [tɪ́ɛ̀], the vowel which appears in the singular is different from the vowel in the plural stem [tɪ́]. There are also vowel alternations in the CD forms, but of a different sort. `Hand' has an extra vowel in the plural rather than in the singular. Also there is a vowel quality difference between the singulars and plurals of `tree', `thief', and `goat'. 
Another noteworthy fact is that the words `eye', `finger', and `toe' are in different declension classes in the two dialects; while they are in the -rɪ plural class in ND, they are in the -r singular class in CD.

\begin{table}

  \caption{-r(ɪ́) plurals in Northern Dagaare and corresponding nouns in Central Dagaare }
  \label{tab:Caragine:rPluralsChange}
  
\begin{tabular}{lllll}
\lsptoprule 
\multicolumn{2}{c}{Northern Dagaare} & \multicolumn{1}{c}{}& \multicolumn{2}{c}{Central Dagaare}\\
Singular & Plural & Gloss 
& Singular & Plural \\
\midrule\midrule
nímíè & nímí-r & `eye' & nímí-rì & *nimi-e\\
nóbíè & nóbí-r & `finger' & núbì-rì & *nubi-e/*nubii-ri\\
g͡bébíè & g͡bébí-r & `toe' & gbɛ́bì-rí & *gbɛbi-e\\
tɪ́ɛ̀ & tɪ́-ɾɪ́ & `tree' & tɪ̀ɛ́ & tìì-rí \citep{bodomoStructure1997}\\
ɲã̀ɲúé & ɲã̀ɲú-r & `thief' & nànɲìgé & *nanɲigi-ri\\
bùó & bó-r & `goat' &bʊ́ɔ́& bʊ́ʊ́-rɪ́ \\
núú & nú-r & `hand' &nû& *nuu-ri\\
\lspbottomrule 

\end{tabular}
\end{table}


%%%%%%%%%%%%%%%%%%%%%%%%%%%%%%%%%%%%%%%%%%%%%%%%%%%
\subsection{-bɛ plurals} \label{sec:Caragine:bePl}

In the -bɛ plural declension class, the singular is unmarked and the plural is marked by -bɛ. The -bɛ plurals surface in two ways: -βɛ and -mɛ. We see -βɛ appear when the suffix is added to plosives final stems and -mɛ appears when it is added to nasal final stems. These appear to be the only stem types which take the -βɛ plural suffix.

As can be seen in \tabref{tab:Caragine:bePluralsCons}, when -bɛ is added to a plosive-final stem, the plosive deletes. Additionally, there is a language general spirantization process which turns [b] into [β] between most sonorants. 

\begin{table}

  \caption{-bɛ plurals in Northern Dagaare and corresponding nouns in Central Dagaare. Central Dagaare data from \citet{aliDictionary2021}.}
  \label{tab:Caragine:bePluralsCons}
  
\begin{tabular}{lllll}
\lsptoprule 
\multicolumn{2}{c}{Northern Dagaare} & \multicolumn{1}{c}{}& \multicolumn{2}{c}{Central Dagaare}\\
Singular & Plural & Gloss 
& Singular & Plural \\
\midrule\midrule
dɛ́b̚ & dɛ̀ɛ̀-βɛ̀ & `man' &dɔ́ɔ́& *dɔ-bɔ\\
jɛ̀b̚ & jɛ̀ɛ̀-βɛ́ & `sibling' &jɔ́ɔ̀& jɔɔmine/*jɔɔre\\
pɔ́g̚ & pɔ́ɔ́-βɛ̀ & `woman' & pɔ́g-ɔ́ & pɔ́g-bɔ́\\
\lspbottomrule 

\end{tabular}
\end{table}

As shown in \tabref{tab:Caragine:bePluralsCons}, there are vowel length alternations between the ND singulars and plurals in the -bɛ class.\footnote{One may note that all of the words in \tabref{tab:Caragine:bePluralsCons} contain -\textsc{atr} mid vowels and end in voiced consonants, but I assume this is a coincidence of the small dataset.} The stem vowel is always short in the singular but long in the plural. However, this alternation is not universal in the class. There is a group of nouns in the -bɛ plural declension class which does not show any vowel alternations (\tabref{tab:Caragine:bePluralNasal}).

\begin{table}

  \caption{-mɛ plurals in Northern Dagaare and corresponding nouns in Central Dagaare. Central Dagaare data from \cite{aliDictionary2021}.}
  \label{tab:Caragine:bePluralNasal}
  
\begin{tabular}{lllll}
\lsptoprule 
\multicolumn{2}{c}{Northern Dagaare} & \multicolumn{1}{c}{}& \multicolumn{2}{c}{Central Dagaare}\\
Singular & Plural & Gloss 
& Singular & Plural \\
\midrule\midrule
nɛ́n & nɛ́mɛ́ & `meat' &nɛ́nɪ̀ & nɛ́mɛ̀ \citep{bodomoStructure1997}\\
bʊ́n & bʊ̀mɛ́ & `thing' &bʊ́nɪ́& *bʊma\\
pán & pàmɛ́ & `door' &pánɪ́ & *pama (`doorway')\\
dàt͡ʃín&  dàt͡ʃímé & `wall' & dànʧíní & *dant͡ʃime \\
k͡pã̀k͡pán & k͡pã̀k͡pámɛ́ & `arm' &kpáŋkpánɪ́& *kpaŋkpama \\
zʊ́m & zʊ́mɛ̀ & `fish' &zʊ́mmʊ́ & *zʊma \\
pìpán & pìpámé & `rock' &píírí& *pie \\
gán& gámɛ́ & `book' &gánɪ̀&gámá \\
sáán & sáámɛ́ & `visitor' & sáánà & saama
\citep{bodomoStructure1997}\\

\lspbottomrule 

\end{tabular}
\end{table}


As alluded to earlier, all of the nouns which end in a nasal consonant in the singular receive the -mɛ suffix in the plural. The -mɛ plural has historically been treated as a separate class \citep{bodomo2007}. However, I argue it is an allomorph of the –bɛ plural when it is added to nasal final stems. When -bɛ́ is added to a nasal-final stem, the nasal and the voiced bilabial plosive [b] coalesce to form a bilabial nasal [m]. The nouns which have a nasal final consonant and take the -bɛ plural suffix do not undergo any sort of vowel alternation. The short vowels in the singular remain short when the plural suffix is added.


%%%%%%%%%%%%%%%%%%%%%%%%%%%%%%%%%%%%%%%%%%%%%%%%%%%%%%%%%%%%%%%%
\section{Towards an analysis} \label{sec:caragine:analysis}

As shown in \sectref{sec:caragine:numberMarking}, several Northern Dagaare nouns exhibit alternations in vowel length and diphthongization between the singular and plural. Seeing data like this in a dialect of Central Dagaare spoken around the towns Jirapa and Ullo,
\citet{anttilaProsodic2009} formed an account for number marking based on minimal prosodic word requirements. In this section, I attempt to apply this account to our Northern Dagaare data, and analyze the instances in which it can capture the variation in our data and the places in which it fails. When their account cannot accurately capture the ND facts, I put forth an alternative analysis. Like \citet{anttilaProsodic2009}, throughout the course of this analysis, I take \textsc{atr} harmony for granted and leave out tone for simplicity.

The section begins with a discussion of -r singulars in \ref{sec:caragine:epluralanal} which function identically to those in Central Dagaare. In \ref{sec:caragine:rianalysis} I show that the word minimallity requirement in CD may not be active in ND and propose an alternative analysis.  \sectref{sec:caragine:bepluralanalysis} discusses the -bɛ plurals, not analyzed in \citet{anttilaProsodic2009, anttilaMetrically2019}, and provides a potential direction for analysis of the class.

%%%%%%%%%%%%%%%%%%%%%%%%%%%%%%%%%%%%%%%%%%%%%%
\subsection{-r singulars} \label{sec:caragine:epluralanal}

As shown earlier, in the -r singular declension class, nouns receive an -r suffix in the singular and an -ɛ suffix in the plural. The plural suffix often results in a dipthong or a long mid vowel. As motivation for a processes which will be discussed later, \citet{anttilaProsodic2009} posit a constraint against derived long mid vowels. However, since -ɛ is a suffix and not an epenthesized vowel, the long mid vowel is allowed. They justify this in OT by having an \textsc{Ident}(V) constraint ranked higher than *VV(mid). These constraints, and the widespread \textsc{Max} constraint are defined below in (\ref{cons}). 

\ea \label{cons} \begin{xlist}
    \ex \textsc{Max}: Assign one violation mark for each segment present in the input which does not have a correspondent in the output.
    \ex  *VV(mid): Assign one violation mark for each long mid vowel.
    \ex \textsc{Ident}(V): For a vowel in the input V\textsubscript{i} which has a corresponding output vowel V\textsubscript{o}, assign one violation for each feature value which is not shared between V\textsubscript{i} and V\textsubscript{o}.
\end{xlist}
\z

Using these constraints, the derivation of the Northern Dagaare word `legs' can be seen in \tabref{tab:Caragine:legs} below. Thus far, this analysis seems to directly capture what we see in Northern Dagaare. The plural suffix -ɛ is added to the stem, resulting in a long mid vowel. Additionally, as the /ɛ/ is in the underlying form, the long mid vowel is allowed to surface.

\begin{table}

  \caption{Tableau deriving [g͡bɛ́ɛ́] `legs' in ND}
  \label{tab:Caragine:legs}
  
\begin{tableau}{c:c|c} 
\inp{/g͡bɛ-ɛ/} 
\const{\textsc{Ident}(V)}
\const{\textsc{Max}}
\const{*VV(mid)}
\cand [\Optimal]{\normalfont{g͡bɛɛ}}  \vio{} \vio{}  \vio{*}
\cand{\normalfont{g͡bɪɛ}} \vio{*!} \vio{}  \vio{}  
\cand{\normalfont{g͡bɛɪ}} \vio{*!} \vio{}  \vio{}  
\cand{\normalfont{g͡bɛ}} \vio{} \vio{*!}  \vio{}  
\end{tableau}
\end{table}

%%%%%%%%%%%%%%%%%%%%%%%%%%%%%%%%%%%%%%%%%%%%%%%%%%
\subsection{-r(ɪ) plurals} \label{sec:caragine:rianalysis}
\largerpage
In the -r(ɪ) plural declension class, the plural receives either the suffix -r or -ɾɪ while the singular is unmarked. As described in \sectref{sec:caragine:rPlurals}, some nouns have a long vowel or diphthong in the plural which is not present in the singular. 

One could posit that one vowel in a VV sequence is deleted when the plural morpheme -r(ɪ) is added, and the lack of vowel shortening in \tabref{tab:Caragine:rPluralsSame} (repeated below in \tabref{tab:Caragine:rPluralsComb} for convenience) is due to a preference for low vowels, as we do not see any low vowels deleting. However, [nàfáʊ̀] ends in a high vowel which seems like it would be a good candidate to delete, yet it remains [nàfáʊ̀ɾɪ́] in the plural. Thus, I follow A\&B in assuming that this is an epenthesis process and the fact that only low vowels are long in the plural is merely an idiosyncrasy of this limited data set. This conclusion is further supported by considering the unpredictability of which vowel appears in the plural form. For example, when the singular form contains a diphthong, [tɪ́ɛ̀] and [bùó], sometimes the first vowel appears in the plural [tɪ́ɾɪ́] and sometimes the second does [bór]. A deletion account could not explain this alternation.

\begin{table}

  \caption{-r(ɪ) plurals}
  \label{tab:Caragine:rPluralsComb}
  
\begin{tabular}{lll}
\lsptoprule 
Singular & Plural & Gloss\\
\midrule\midrule
báá & báá-r & `dog'\\
bàà & bàà-r & `river'\\
dàà & dáá-r & `stick'\\
nàfáʊ̀ & nàfáʊ̀-ɾɪ́ & `shoe'\\
gàdò & gàdò-ɾí & `bed'\\
\midrule\midrule
nímíè & nímí-r & `eye'\\
nóbíè & nóbí-r & `finger'\\
g͡bébíè & g͡bébí-r & `toe'\\
tɪ́ɛ̀ & tɪ́-ɾɪ́ & `tree'\\
ɲã̀ɲúé & ɲã̀ɲú-r & `thief'\\
bùó & bó-r & `goat'\\
núú & nú-r & `hand'\\
\lspbottomrule 

\end{tabular}
\end{table}

\citet{anttilaProsodic2009} state that in Central Dagaare, when a stem is monomoraic, a vowel is epenthesized in the singular to create a fully formed bimoraic foot. They state that this is a constraint seen across the language; all phonological words must minimally contain a fully formed prosodic foot. To satisfy this constraint, a vowel is added to the end of monomoraic stems. This vowel matches the quality of the previous vowel. However, as can be seen in \tabref{tab:Caragine:rPluralsComb}, the vowels do not always match. Anttila and Bodomo argue that this is due to language general constraints against word final high vowels and long mid vowels (briefly mentioned above).

The following tableau are modeled after \cite{anttilaProsodic2009}, but the data are from Northern Dagaare. The relevant constraints they describe are as follows: 

\ea \label{ex:Caragine:HeIsInsideTheHouse} \begin{itemize}
    \item[a.] \textsc{Dep}(V):	Assign one violation mark for each vowel in the output which does not have a correspondent in the input.\\
    \item[b.] \textsc{FtBin}: Assign one violation mark for each foot which is not minimally bimoraic.\\
    \item[c.] *I]:	Assign one violation mark for each word-final [+high] vowel.\\
\end{itemize}
\z

\begin{table}

  \caption{Tableau deriving [tɪɛ] `tree' in ND}
  \label{tab:Caragine:tree}

\begin{tableau}{c|c:c} 
\inp{/tɪ/} 
\const{\textsc{FtBin}} 
\const{\textsc{Dep}(V)} 
\const{*I]} 
\cand{\normalfont{tɪ}} \vio{*!} \vio{} \vio{*} 
\cand{\normalfont{tɪɪ}} \vio{} \vio{*} \vio{*!}  
\cand[\Optimal]{\normalfont{tɪɛ}} \vio{} \vio{*} \vio{}
\end{tableau}
\end{table}


\tabref{tab:Caragine:tree} shows the derivation of [tɪ́ɛ̀] ‘tree’. The stem is monomoraic, so a vowel is epenthesized to satisfy \textsc{FtBin}. You would normally epenthesize a vowel with the same features as the stem vowel. However, to satisfy the constraint against word final high vowels, a mid vowel is added instead. 

But, as can be seen in \tabref{tab:Caragine:rPluralsComb}, there are several words which end in high vowels. Similar to the analysis in \ref{sec:caragine:epluralanal} of the *VV(mid) constraint, the constraint against word final high vowels only affects ``derived vowels'' -- i.e., those which do not have a corespondent in the UR. This is again attained by a higher ranked \textsc{Ident}(V) which prevents vowel lowering if the violating high vowel has a correspondent in the input. A simple case of this can be seen in \tabref{tab:Caragine:bed}.

\begin{table}

  \caption{Tableau deriving [gadoɾi] `bed' in Northern Dagaare}
  \label{tab:Caragine:bed}

\begin{tableau}{c:c|c} 
\inp{/gado-ɾi/}
\const{\textsc{Ident(V)}} 
\const{\textsc{Dep}(V)}
\const{*I]} 
\cand[\Optimal]{\normalfont{gadoɾi}} \vio{} \vio{}\vio{*} 
\cand{\normalfont{gadoɾe}} \vio{*!} \vio{}\vio{} 
\cand{\normalfont{gadoɾie}} \vio{} \vio{*!}\vio{}
\end{tableau}
\end{table}

\noindent Despite violating *I], the [i] remains in the output because changing it would violate the higher ranked \textsc{Ident(V)}. 

As noted in \sectref{sec:caragine:rPlurals}, the Central Dagaare plural suffix, -rɪ, often appears without the vowel in Northern Dagaare. This is important because while the plural marker straight forwardly satisfies the \textsc{FtBin} constraint in Central Dagaare, in Northern Dagaare it is necessary to say that /r/ is moraic so as to satisfy the foot minimality requirement and not trigger vowel epenthesis in the plural. The moraic status of coda consonants will be discussed further in  \sectref{sec:caragine:bepluralanalysis}.


\begin{table}

  \caption{Tableau deriving [buo] `goat' in Northern Dagaare}
  \label{tab:Caragine:goat}

\begin{tableau}{c|c|c:c} 
\inp{\ips{bo}} 
\const{\textsc{FtBin}} 
\const{\textsc{Dep}(V)} 
\const{*I]} 
\const{*VV(mid)}
\cand{bo} \vio{*!} \vio{} \vio{} \vio{} 
\cand{boo} \vio{} \vio{*} \vio{} \vio{*!} 
\cand[\Optimal]{buo} \vio{} \vio{*} \vio{} \vio{} 
\cand{bou} \vio{} \vio{*} \vio{*!} \vio{} 
\end{tableau}
\end{table}

In \tabref{tab:Caragine:goat}, we see the derivation of [bùó] ‘goat’. Again, the stem is mono-moraic so a vowel must be epenthesized. The constraint against long mid vowels prevents the normal epenthesis of [o] and instead a high vowel is epenthesized. However, due to the constraint against word final high vowels, the epenthesized vowel cannot occur word finally as it has in the other examples. Rather, the  high vowel [u] is epenthesized before the stem vowel.\footnote{One may note that there is nothing particular indicating where the epenthesized vowel appears or what quality it should take. This is a shortcoming of the original analysis. Since I will soon argue that this analysis does not work for ND and propose an alternative, I will not try to plug this hole. However, future work on CD would benefit from a more complete analysis of the nature and location of the epenthesized vowel.}

This accounts well for Bodomo and Anttila’s data as well as the Northern Dagaare data discussed above, but when attempting to expand this theory to a larger set of ND nouns, it does not capture all the variation. The word minimality account does not explain the epenthetic vowel in the singulars of nouns with multi-moraic stems (\tabref{tab:Caragine:epenth}).

\begin{table}

  \caption{Multimoraic Stems of ND -rɪ plurals}
  \label{tab:Caragine:epenth}
  
\begin{tabular}{llll}
\lsptoprule 
Stem & Singular & Plural & Gloss\\
\midrule\midrule
nóbì & nóbíé  & nóbí-r &  `finger' \\
nímí &  nímíè &  nímí-r &  `eye'\\
g͡bébí& g͡bébíè & g͡bébí-r & `toe'\\
tèɾét͡ʃì  & tèɾét͡ʃìé  & tèɾét͡ʃì-ɾí &  `train'\\
ɲã̀ã́ɲù  &  ɲã̀ã́ɲùé &  ɲã̀ã́ɲù-ɾí &  `thief'\\
\lspbottomrule 

\end{tabular}
\end{table}


As the stems in \tabref{tab:Caragine:epenth} are already, at minimum, bimoraic, Anttila and Bodomo’s account does not provide an explanation for the epenthetic vowels in the singular. In fact, as seen in \tabref{tab:Caragine:finger}, their constraints would make the wrong prediction. 

\begin{table}
  \caption{Tableau of the failed derivation of [nóbíé] `finger' in ND}
  \label{tab:Caragine:finger}

\begin{tableau}{c|c|c:c} 
\inp{\ips{nobi}} 
\const{\textsc{FtBin}} 
\const{\textsc{Dep}(V)} 
\const{*I]} 
\const{*VV(mid)}
\cand[\bomb]{nobi} \vio{} \vio{} \vio{*} \vio{} 
\cand[\grimace]{nobie} \vio{} \vio{*!} \vio{} \vio{} 
\cand{nobii} \vio{} \vio{*!} \vio{*} \vio{} 
\end{tableau}
\end{table}

Finally, the prosodic account put forth by Anttila and Bodomo is dependent on the fact that Dagaare requires a fully formed foot for all prosodic words. They acknowledge that there are a few exceptions to this rule for Central Dagaare, specifically [bâ] ‘father’, [mǎ] ‘mother’, [nû] ‘hand’, [zû] ‘head’, all of which are very high frequency words. However, deficient feet seem to be much more common in Northern Dagaare. In addition to those mentioned for Central Dagaare, the Northern dialect also has [nú] ‘face,’ as well as many verbs including but not limited to [sʊ̀] ‘bathe’, [ɲʊ̀] ‘drink’, and [tʃɛ́] ‘slice.’ 

There may  be a concern that the verbs do not in and of themselves form full prosodic units as they typically exist as part of a verbal paradigm including auxiliaries, tense markers, and object pronouns. However, we see that verbs can stand alone in imperatives \REF{ex:Caragine:come}.

\ea \label{ex:Caragine:come}
wá \\
‘come’
\z

Word minimality does not seem to be as active of a constraint in Northern Dagaare as it was in Central Dagaare. Thus, using it as a motivation for vowel epenthesis in nominals is less obvious. Additionally, epenthesis occurs in some words which already satisfy the proposed word minimality.

Rather, following work on quantity-manipulating morphology \citep{fujimura, trommer}, I propose that there is an empty mora suffixed to all singular nouns in the -rɪ plural class. This mora is then filled by linking to the adjacent vowel. I will generally assume an autosegmental feature geometry such as those proposed in \citet{Clements}, but the specifics of the system are not central to the analysis. Having a mora be the affix which spells out singular number requires that all nouns in this class will have an ``epenthetic'' vowel due to \textsc{Max}-$\mu$. This must be ranked over \textsc{1-to-1}, a constraint against lengthening/diphthongization of the vowel.

\ea \begin{xlist}
    \ex \textsc{Max}-$\mu$: assign one violation mark for each mora in the input which is not filled in the output. 
    \ex  \textsc{1-to-1}: Assign one violation mark for each segment  that is linked to more than one prosodic node.
    \end{xlist}
\z

The same constraints proposed by \cite{anttilaProsodic2009} can be used to account for the instances in which the vowel qualities do not match. If after linking your final vowel to the empty mora, the final vowel is high, the vowel is lowered to satisfy *I] (\tabref{tab:Caragine:morafinger}). If after linking, there is a long mid vowel, in violation of *VV(mid), the first element of the vowel sequence is raised (\tabref{tab:Caragine:moragoat}).

\begin{table}
  \caption{Tableau illustrating high vowel lowering in [nóbíé] `finger' in ND}
  \label{tab:Caragine:morafinger}
\begin{tableau}{c:c|c}
\inp{\ips{\normalfont{nobi\textsuperscript{$\mu-\mu$}}}}  
\const{\textsc{Max}-$\mu$}
\const{*I]} 
\const{\textsc{1-to-1}}
\cand{\normalfont{nobi\textsuperscript{$\mu$}}} \vio{*!} \vio{*} \vio{} 
\cand[\Optimal]{\normalfont{nobi\textsuperscript{$\mu$}e\textsuperscript{$\mu$}}} \vio{} \vio{} \vio{}
\cand{\normalfont{nobi\textsuperscript{$\mu$$\mu$}}} \vio{} \vio{*!} \vio{*}
\end{tableau}
\end{table}

\begin{table}
  \caption{Tableau illustrating mid vowel raising in [buo] `goat' in ND}
  \label{tab:Caragine:moragoat}
\begin{tableau}{c:c|c}
\inp{\ips{bo\textsuperscript{$\mu-\mu$}}}  
\const{\textsc{Max}-$\mu$}
\const{*VV(mid)}
\const{\textsc{1-to-1}}
\cand{\normalfont{bo\textsuperscript{$\mu$}}} \vio{*!} \vio{} \vio{} 
\cand[\Optimal]{\normalfont{bu\textsuperscript{$\mu$}o\textsuperscript{$\mu$}}} \vio{} \vio{} \vio{}
\cand{\normalfont{bo\textsuperscript{$\mu$$\mu$}}} \vio{} \vio{*!} \vio{*}
\end{tableau}
\end{table}

As shown in \tabref{tab:Caragine:morastick}, not all singular nouns in this class undergo vowel lengthening. There is no lengthening in stems which already contain a long vowel. This can easily be explained, as there are no ``superlong'' vowels in Dagaare. 

\begin{table}
  \caption{Tableau showing no lengthening of [daa] `stick' in ND}
  \label{tab:Caragine:morastick}
\begin{tableau}{c|c|c|c}
\inp{\ips{da\textsuperscript{$\mu$$\mu$}-$\mu$}}  
\const{*\textsc{SuperLong}}
\const{\textsc{Max}-$\mu$}
\cand{\normalfont{da\textsuperscript{$\mu$$\mu$$\mu$}}} \vio{*!} \vio{}
\cand[\Optimal]{\normalfont{da\textsuperscript{$\mu$$\mu$}}} \vio{} \vio{*}
\end{tableau}
\end{table}

Under this analysis, in the -r(ɪ) plural class, nouns in the singular -- so long as they do not already end in a long vowel -- will always undergo vowel lengthening or diphthongization due to the added mora. This contrasts with Anttila and Bodomo's word minimallity analysis which would only predict vowel lengthening/diphthongization on stems that are less than a fully formed foot.

It should be noted that this analysis would not correctly predict the ND singulars [gàdò] `bed' and [núú] `hand' . Mora affixation accompanied by mid vowel raising would predict [gàdùò]. Additionally, the constraint *I] would predict high vowel lowering in `hand' to produce [núó]. However, this would also be incorrectly predicted under Anttila and Bodomo's account.\footnote{In CD [nû] `hand' is also a difficult case which would not be predicted under Anttila and Bodomo's word minimality account.}

%%%%%%%%%%%%%%%%%%%%%%%%%%%%%%%%%%%%%%%%%%%%%%%%%%%%%%%%%%%%%%%
\subsection{-bɛ́ plurals} \label{sec:caragine:bepluralanalysis}

Anttila and Bodomo did not address the -bɛ plurals in their work on Central Dagaare.\footnote{This in unsurprising since, as can be seen in Tables \ref{tab:Caragine:bePluralsCons} and \ref{tab:Caragine:bePluralNasal}, it is likely that the -bɛ plurals do not actually form a unified class in CD.} As shown in \tabref{tab:Caragine:bePluralsCons} the -bɛ́ plurals in ND have a short vowel in the singular and a long vowel in the plural. Logically, this could be due to shortening in the presence of a coda or lengthening with the deletion of an underlying coda. However, there are many words with long vowels and codas (all of the -r plurals with long vowels as well as several other words including [wáàb] `snake', [bɛ́ɛ́n] `one', and [lɛ̀ɛ̀m] `taste') so it seems unlikely that the vowel is shortened in the presence of a coda.

Therefore, I conclude that the -bɛ plurals are undergoing a process of compensatory lengthening, in which one segment is lengthened to compensate for the loss of another. For example, /pɔ̀g-bɛ́/ `women' becomes [pɔ́ɔ́βɛ̀]. Here the vowel lengthens to compensate for the loss of the [g]. Additionally, there is a language general process of /b/ spirantization between sonorants accounted for by a \textsc{Spirant} constraint. Compensatory lengthening cases such as these are frequently analyzed with a moraic account (\citealp{hayes1989, shaw}). A moraic coda, when deleted, leads to an empty mora which is then filled by lengthening of the preceding vowel. I consider this deletion to be triggered by a *CC constraint specified to ban adjacent [-sonorant] consonants. I propose this holds across the language, as we never see two adjacent obstruents. Additionally, the deletion of the first rather then the second obstruent is due to the ranking of \textsc{Max}$_{onset}$ over a general \textsc{Max} constraint.

\ea \begin{xlist}
    \ex *CC: Assign one violation mark for each consonant cluster in which both segments are [-sonorant].
    \ex\textsc{Max}$_{onset}$: Assign one violation mark for each onset in the input which does not have a corespondent  in the output.
    \ex \textsc{Spirant}: Assign one violation mark for each [-son] segment between two [+son] segments.
\end{xlist}
\z

\begin{table}

  \caption{Tableau deriving [pɔ̀ɔ̀βɛ́] `women'}
  \label{tab:Caragine:women}
  
\begin{tableau}{c:c:c:c|c:c} 
\inp{\ips{\normalfont{pɔ\textsuperscript{$\mu$}g\textsuperscript{$\mu$}-bɛ}}}
\const{*CC}
\const{\textsc{Max}$_{onset}$}
\const{\textsc{Max}-$\mu$}
\const{\textsc{Spirant}}
\const{\textsc{1-to-1}}
\const{\textsc{Max}}
\cand {\normalfont{pɔ\textsuperscript{$\mu$}g\textsuperscript{$\mu$}bɛ}}  \vio{*!} \vio{}  \vio{} \vio{}\vio{}\vio{}
\cand{\normalfont{pɔ\textsuperscript{$\mu$$\mu$}bɛ}} \vio{} \vio{}  \vio{} \vio{*!} \vio{*}\vio{*}
\cand{\normalfont{pɔ\textsuperscript{$\mu$$\mu$}ɣɛ}} \vio{} \vio{*!}  \vio{} \vio{} \vio{*}\vio{*}
\cand[\Optimal]{\normalfont{pɔ\textsuperscript{$\mu$$\mu$}βɛ}} \vio{} \vio{}  \vio{}  \vio{}\vio{*}\vio{*}
\cand{\normalfont{pɔ\textsuperscript{$\mu$}βɛ}} \vio{} \vio{}  \vio{*!}  \vio{}\vio{}\vio{*}
\end{tableau}
\end{table}

In \tabref{tab:Caragine:women} *CC in conjunction with \textsc{Max}$_{onset}$ causes the deletion of the coda [g]. In compensation for this, triggered by \textsc{Max}-$\mu$, the vowel lengthens to fill the empty mora. 

However, as seen in \sectref{sec:Caragine:bePl}, when the -bɛ plural suffix attaches to nasal final roots, the /b/ and the final nasal coalesce into an [m] -- [pán] `door', [pàmɛ́] `doors.' Thus, \textsc{Uniformity}, a constraint preventing coalescence, must be ranked below *NC, which triggers coalescence. Additionally, since the [nasal] and [labial] features both surface on the resulting segment when the [n] and [b] coalesce, I introduce \textsc{Max}$_{[nasal]}$ and \textsc{Max}$_{[lab]}$.

\ea \begin{xlist}
    \ex \textsc{Uniformity}: assign one violation mark for each element in the output that has more than one correspondent in the input.
    \ex \textsc{Max}\textsubscript{[nasal]}: assign one violation mark for each [nasal] feature in the input that has no correspondent in the output.
    \ex \textsc{Max}\textsubscript{[lab]}: assign one violation mark for each [labial] feature in the input that has no correspondent in the output.
    \ex *NC: assign one violation mark for each nasal-consonant sequence.
\end{xlist}
\z


\begin{table}

  \caption{Tableau illustrating the failed derivation of [pàmɛ́] `door'}
  \label{tab:Caragine:doors}
\fittable{
\begin{tableau}{c:c:c:c:c|c:c:c} 
\inp{\ips{\normalfont{pa\textsuperscript{$\mu$}n$_{1}^\mu$-b$_{2}$ɛ}}}
\const{\textsc{Max}\textsubscript{[lab]}}
\const{\textsc{Max}\textsubscript{${onset}$}}
\const{*NC}
\const{\textsc{Max}\textsubscript{[nasal]}}
\const{\textsc{Max}-$\mu$}
\const{\textsc{1-to-1}}
\const{\textsc{Max}}
\const{\textsc{Uni}}
\cand {\normalfont{pa\textsuperscript{$\mu$}n$_{1}^\mu$.b$_{2}$ɛ}}  \vio{}\vio{}\vio{*!} \vio{}  \vio{} \vio{}\vio{}\vio{}
\cand{\normalfont{pa\textsuperscript{$\mu$$\mu$}.b$_{2}$ɛ}} \vio{}\vio{}\vio{} \vio{*!}  \vio{} \vio{*}\vio{*} \vio{}
\cand{\normalfont{pa\textsuperscript{$\mu$}.b$_{2}$ɛ}}\vio{} \vio{} \vio{} \vio{*!}  \vio{*}\vio{} \vio{*} \vio{}
\cand{\normalfont{pa\textsuperscript{$\mu$$\mu$}.n$_{1}$ɛ}}  \vio{*!}\vio{*}\vio{} \vio{}  \vio{} \vio{*} \vio{*}\vio{}
\cand[\grimace]{\normalfont{pa\textsuperscript{$\mu$}.m$_{12}$ɛ}} \vio{}\vio{} \vio{} \vio{}  \vio{*!}\vio{} \vio{} \vio{*}
\cand[\bomb]{\normalfont{pa\textsuperscript{$\mu$$\mu$}.m$_{12}$ɛ}} \vio{}\vio{} \vio{} \vio{}  \vio{} \vio{*}\vio{} \vio{*}
\cand{\normalfont{pa\textsuperscript{$\mu$$\mu$}.n$_{12}$ɛ}} \vio{*!}\vio{}\vio{} \vio{}  \vio{}\vio{*} \vio{} \vio{*}
\end{tableau}
}
\end{table}

When a nasal and [b] coalesce, there is no change in vowel length. As can be seen in \tabref{tab:Caragine:doors}, the same ranking which correctly captured compensatory lengthening in obstruent final roots fails to correctly predict the behavior of nasal final roots. If we assume Onset Maximization \citep{Ito, Selkirk} the [m] would be the onset of the second syllable, not the coda of the first. Therefore it would not be able to hold a mora, and [paamɛ] would incorrectly win to satisfy \textsc{Max}-$\mu$. In some way, the retention of the nasal feature from the underlying coda must ``count'' as filling a second mora. One potential explanation for this would be that the resulting [m] is ambisyllabic; this way it could satisfy Onset Maximization but still be able to hold a mora. \tabref{tab:Caragine:doors-correct} shows what this could look like in an OT tableau.

\begin{table}

  \caption{Tableau deriving [pàmɛ́] `door'}
  \label{tab:Caragine:doors-correct}
  
\begin{tableau}{c:c|c:c:c} 
\inp{\ips{\normalfont{pan\textsubscript{1}-b\textsubscript{2}ɛ}}}
\const{*NC}
\const{\textsc{Max}-$\mu$}
\const{\textsc{1-to-1}}
\const{\textsc{Max}}
\const{\textsc{Uni}}
\cand{\normalfont{pa\textsuperscript{$\mu$}.m\textsubscript{12}ɛ}}  \vio{} \vio{*!}  \vio{} \vio{} \vio{*}
\cand[\Optimal]{\normalfont{pa\textsuperscript{$\mu$}(m$^\mu_{12}$)ɛ}} \vio{}  \vio{}  \vio{} \vio{} \vio{*}
\cand{\normalfont{pa\textsuperscript{$\mu\mu$}.m$_{12}$ɛ}} \vio{}\vio{}\vio{*} \vio{} \vio{*}
\end{tableau}
\end{table}

In \tabref{tab:Caragine:doors-correct}, Candidates (a) and (b) are segmentally identical, but in (b) one of the mora in the first syllable is filled by some of the feature information of the underlying /n/. If the [m] is ambisyllabic, filling both the coda of the first syllable and the onset of the second, we may expect to see some phonetic evidence. I leave investigation of this possibility to future research.

%%%%%%%%%%%%%%%%%%%%%%%%%%%%%%%%%%%%%%%%%%%%%%%%%%%%%%
\section{Conclusion} \label{sec:caragine:conclusion}

\citet{anttilaProsodic2009} formulated an analysis of vowel alternations in Central Dagaare. Their analysis of the -r singulars seems to map directly to the -r singulars in Northern Dagaare. However, they argue that for Central Dagaare, vowels are epenthesized in the singular to fulfill a word minimality requirement for bimoraic feet. But, I have shown that in Northern Dagaare some words have epenthesized vowels even when they already contain bimoraic feet. Instead, I propose that the singular marker for the -r(ɪ) plural class is an empty mora. This mora suffix is then filled by linking to the adjacent vowel, and language general constraints formulated by Anttila and Bodomo modify the vowels to form the correct surface forms.

Additionally, I put forth a preliminary account for vowel length alternations in the -βɛ and -mɛ plurals, here analyzed as a single class, explaining why compensatory lengthening occurs in the former but not the latter. While more data is necessary to help clarify these patterns and formulate a more comprehensive analysis, I have provided a stepping off point for which to analyze the vowel alternations in Northern Dagaare number marking. 

%%%%%%%%%%%%%%%%%%%%%%%%%%%%%%%%%%
\section*{Acknowledgments}
I would like to extend my thanks to Michael Obiri-Yeboah, Kate Mooney, and two anonymous reviewers for their assistance formulating this paper. Additionally, I thank the conference attendees of ACAL 54 as well as the University of Maryland Phon-Circle for providing insightful avenues of investigation. Finally, I would like to thank the Georgetown University Field Methods class of 2022 and our language informant, without whom this work would not be possible.

\sloppy %this command eliminates a quirk that appears in the bibliography, you can just leave it here.

\printbibliography[heading=subbibliography,notkeyword=this]

\end{document}
